% --- Définition des couleurs institutionnelles (Bleu Nuit & Bleu Clair) ---
% Couleurs historiques des tableaux du chapitre 4.
\colorlet{BleuNuit}{argostablehead}
\colorlet{BleuClair}{argostablerow}
% La couleur des filets est gérée globalement par le préambule.
% Certaines lignes historiques appliquent \rowcolor dans chaque cellule.
% \cellcolor conserve le rendu attendu sans provoquer d'erreur d'alignement.
\renewcommand{\rowcolor}[1]{\cellcolor{#1}}

\chapter{Étude de la viabilité}
\label{chap:viabilite_financière}

L’étude de la viabilité financière permet d’évaluer la rentabilité prévisionnelle pour notre projet ARGOS a généré une technique de gestion qui a pour objectif de donner une vision synthétique de la situation financière d’un projet. Elle consiste également à valoriser les flux financiers résultant des études précédentes afin de déterminer la rentabilité d’un projet, tout ceci sur la base des prévisions financières qui seront ressorties. Cette étude constitue un outil d’aide à la décision stratégique. Elle s’adresse aux porteurs de projet, aux partenaires financiers et aux investisseurs désireux de s’assurer de la solidité économique du projet. Pour y parvenir, nous allons élaborer les différents états financiers prévisionnels ensuite présenter quelques indicateurs de la rentabilité du projet.

\section{Prévisions financières}
L’établissement de nos prévisions financières s’appuie sur les différentes hypothèses issues de l’étude de marché et de l’analyse stratégique menées dans les chapitres précédents. Cette partie nous permet d’évaluer le coût total de notre projet et de ressortir l’origine des ressources financières. Pour réaliser notre projet, nous aurons besoin de matériels informatiques, logiciels et matériels de bureau qui favoriseront le démarrage de notre activité. Nous allons aussi évaluer le montant de la charge salariale qui nous aidera à estimer notre besoin en fonds de roulement initial (BFR).

% \\\\\\\\\\\\\\\\\\\\\\\\\\\\\\\\\\\\\\\\\\\\\\\\\\\\\\\\\\\\\  analyse previsionnels


\subsection{Analyse fréquentielle des flux de conteneurs en Afrique centrale}

La rentabilité du modèle économique ARGOS repose sur un nombre suffisant de trajets de conteneurs pour générer les revenus nécessaires à couvrir les coûts d’investissement et d’exploitation. Cette section présente une analyse approfondie des données réelles des principaux ports d’Afrique centrale.
% \\\\\\\\\\\\\\\\\\\\\\\\\  ajouter d'ici 

\subsubsection{Port de Douala (Cameroun) : Pôle historique}
Le Port Autonome de Douala (PAD) constitue le hub maritime le plus important d’Afrique centrale, concentrant entre 71\,\% et 76\,\% du volume global de trafic maritime camerounais selon l’analyse sectorielle du transport maritime camerounais publié le 1er février 2026\cite{pad2025}.

\begin{table}[ht]
\centering
\caption{Trafic du Port de Douala (2023-2024)} \label{tab:trafic_douala}
\begin{tabularx}{\textwidth}{@{}Xrrc@{}}
\toprule
\rowcolor{BleuNuit}\color{white}\textbf{Indicateur} & \rowcolor{BleuNuit}\color{white}\textbf{2023} & \rowcolor{BleuNuit}\color{white}\textbf{2024} & \rowcolor{BleuNuit}\color{white}\textbf{Évolution} \\
\midrule
Traffic export-import (tonnes) & 12 190 000 & 12 920 000 & +6,0\,\% \\
\rowcolor{BleuClair}Mouvements de navires & 1 450 & 1 520 & +4,8\,\% \\
Durée moyenne de service (jours) & 5,2 & 5,2 & — \\
\bottomrule
\end{tabularx}
\end{table}
\tablesourcepromoteurs

L'analyse de l'activité du Port Autonome de Douala met en évidence une dynamique de croissance soutenue de ses plateformes logistiques\cite{pad2025}. 
Entre 2023 et 2024, le trafic global des marchandises à l'exportation et à l'importation a progressé de 6\,\% 
pour atteindre 12,92 millions de tonnes (soit une augmentation notable de plus de 720\,000 tonnes sur une période d'un an)\cite{pad2025}. 
Cette vitalité opérationnelle s'accompagne d'une hausse du trafic maritime, matérialisée par 1\,520 mouvements de navires en 2024 (en augmentation de 4,8\,\% par rapport à l'exercice précédent), 
venant s'inscrire dans une trajectoire globale d'activité cumulée à plus de 5\,200 escales sur la période quinquennale récente\cite{pad2025}. 
Enfin, la durée moyenne de séjour et de service des unités (maintenue de manière constante à 5,2 jours) témoigne d'une stabilisation des opérations de traitement portuaire\cite{pad2025}, 
tout en soulignant la marge d'optimisation temporelle que des solutions de traçabilité avancée comme ARGOS peuvent 
contribuer à combler.

\subsubsection{Port de Kribi (Cameroun) : Croissance exponentielle}
Le Port Autonome de Kribi (PAK) connaît une dynamique de croissance exceptionnelle, avec un trafic conteneurisé multiplié par plus de quatre entre 2018 et 2025\cite{pak2026}. Cette affirmation est confirmée par les rapports officiels et les analyses sectorielles récents sur la performance portuaire du Cameroun, qui soulignent la progression spectaculaire du Port Autonome de Kribi (PAK) entre 2018 et 2025\cite{pak2026}.

\begin{xltabular}{\textwidth}{@{} X r r @{}}
\caption{Évolution du trafic conteneurs à Kribi (2018-2025)} \label{tab:trafic_kribi} \\
\toprule
\rowcolor{BleuNuit}\color{white}\textbf{Année} & \rowcolor{BleuNuit}\color{white}\textbf{Trafic conteneurs (EVP)} & \rowcolor{BleuNuit}\color{white}\textbf{Évolution annuelle} \\
\midrule
\endfirsthead
\multicolumn{3}{c}{\tablename\ \thetable{} -- Suite de la page précédente} \\
\toprule
\rowcolor{BleuNuit}\color{white}\textbf{Année} & \rowcolor{BleuNuit}\color{white}\textbf{Trafic conteneurs (EVP)} & \rowcolor{BleuNuit}\color{white}\textbf{Évolution annuelle} \\
\midrule
\endhead
\midrule
\multicolumn{3}{r}{\textit{Suite à la page suivante...}} \\
\endfoot
\bottomrule
\multicolumn{3}{l}{\textit{Source : Les Promoteurs 2025-2026}} \\
\endlastfoot
2018 & 135 820 & — \\
\rowcolor{BleuClair}2019 & 153 403 & +13,0\,\% \\
2020 & 239 871 & +56,4\,\% \\
\rowcolor{BleuClair}2021 & 210 042 & -12,4\,\% \\
2022 & 277 862 & +32,3\,\% \\
\rowcolor{BleuClair}2023 & 275 063 & -1,0\,\% \\
2024 & 304 546 & +10,7\,\% \\
2025 & 555 398 & +82,4\,\% \\
\end{xltabular}

L'analyse de l'évolution du trafic de conteneurs au Port Autonome de Kribi entre 2018 et 2025 révèle une trajectoire d'expansion remarquable, illustrant la montée en puissance de ce hub portuaire stratégique\cite{pak2026}. Porté par une croissance globale spectaculaire de plus de 308\,\% sur la période (le trafic annuel passant de 135\,820 EVP à 555\,398 EVP) et un taux de croissance annuel moyen s'établissant à 22,3\,\%, le port enregistre une accélération sans précédent lors de l'exercice 2025 avec un bond de 82,4\,\% (portant le tonnage global à 12,7 millions de tonnes)\cite{pak2026}. Cette dynamique positionne Kribi comme un acteur incontournable de la façade maritime nationale, consolidant sa part de marché à 37,5\,\% en 2025\cite{pak2026}. Cette forte densité de flux conteneurisés renforce d'autant plus la pertinence d'implanter des solutions de sécurisation et de télémétrie intelligente telles qu'ARGOS pour optimiser ces flux logistiques majeurs.

\subsubsection{Synthèse des flux conteneurisés en Afrique centrale}

\begin{xltabular}{\textwidth}{@{} X r r r @{}}
\caption{Synthèse des flux conteneurisés par port} \label{tab:synthese_flux} \\
\toprule
\rowcolor{BleuNuit}\color{white}\textbf{Port} & \rowcolor{BleuNuit}\color{white}\textbf{Trafic 2023 (EVP)} & \rowcolor{BleuNuit}\color{white}\textbf{Trafic 2024 (EVP)} & \rowcolor{BleuNuit}\color{white}\textbf{Trafic 2025 (EVP)} \\
\midrule
\endfirsthead
\multicolumn{4}{c}{\tablename\ \thetable{} -- Suite de la page précédente} \\
\toprule
\rowcolor{BleuNuit}\color{white}\textbf{Port} & \rowcolor{BleuNuit}\color{white}\textbf{Trafic 2023 (EVP)} & \rowcolor{BleuNuit}\color{white}\textbf{Trafic 2024 (EVP)} & \rowcolor{BleuNuit}\color{white}\textbf{Trafic 2025 (EVP)} \\
\midrule
\endhead
\midrule
\multicolumn{4}{r}{\textit{Suite à la page suivante...}} \\
\endfoot
\bottomrule
\multicolumn{4}{l}{\textit{Source : Les Promoteurs 2025-2026}} \\
\endlastfoot
Douala & 190 000 & 195 000 & 200 000 \\
\rowcolor{BleuClair}Kribi & 275 063 & 304 546 & 555 398 \\
Pointe-Noire (Congo) & 1 000 000 & 1 050 000 & 1 100 000 \\
\rowcolor{BleuClair}Matadi (RDC) & 200 000 & 215 000 & 230 000 \\
Autres ports & 100 000 & 105 000 & 110 000 \\
\midrule
\textbf{Total Afrique centrale} & \textbf{1 765 063} & \textbf{1 869 546} & \textbf{2 195 398} \\
\end{xltabular}
% \\\\\\\\\\\\\\\\\\\\\\\\\\\ fin ici 

\subsubsection{Estimation du nombre de trajets par dispositif}
Sur la base des données disponibles, nous établissons les hypothèses suivantes :

\begin{xltabular}{\textwidth}{@{} X c X @{}}
\caption{Hypothèses de calcul des trajets} \label{tab:hypotheses_trajets} \\
\toprule
\rowcolor{BleuNuit}\color{white}\textbf{Paramètre} & \rowcolor{BleuNuit}\color{white}\textbf{Valeur} & \rowcolor{BleuNuit}\color{white}\textbf{Justification} \\
\midrule
\endfirsthead
\multicolumn{3}{c}{\tablename\ \thetable{} -- Suite de la page précédente} \\
\toprule
\rowcolor{BleuNuit}\color{white}\textbf{Paramètre} & \rowcolor{BleuNuit}\color{white}\textbf{Valeur} & \rowcolor{BleuNuit}\color{white}\textbf{Justification} \\
\midrule
\endhead
\midrule
\multicolumn{3}{r}{\textit{Suite à la page suivante...}} \\
\endfoot
\bottomrule
\multicolumn{3}{l}{\textit{Source : Les Promoteurs 2025-2026}} \\
\endlastfoot
Durée moyenne de stationnement & 5,2 jours & Données PAD 2025 \\
\rowcolor{BleuClair}Nombre moyen de rotations annuelles & 4 à 6 & Estimation conservatrice \\
Taux d’occupation effectif & 85\,\% & Basé sur les données de trafic \\
\rowcolor{BleuClair}Nombre de trajets par an (conservateur) & 4 & Rotation minimale \\
Nombre de trajets par an (modéré) & 5 & Rotation moyenne \\
Nombre de trajets par an (optimiste) & 6 & Rotation maximale \\
\end{xltabular}

\noindent Le Tableau~\ref{tab:hypotheses_trajets} pose les bases opérationnelles et quantitatives indispensables à l'évaluation du modèle économique d'ARGOS. En s'appuyant sur une durée moyenne de stationnement de 5,2 jours (alignée sur les indicateurs du port de Douala) et en retenant une fourchette de rotations annuelles allant de 4 à 6 trajets, l'analyse intègre un coefficient de réalisme et de prudence (taux d'occupation effectif de 85\,\%). Cette modélisation par paliers au conservateur, modéré et optimiste et permet d'encadrer les projections financières futures et garantit que les calculs de rentabilité et de chiffre d'affaires s'appuieront sur des hypothèses de rotation de flotte rigoureusement maîtrisées.

\subsubsection{Marché adressable}

\begin{xltabular}{\textwidth}{@{} X r c c @{}}
\caption{Marché adressable en Afrique centrale} \label{tab:marche_adressable} \\
\toprule
\rowcolor{BleuNuit}\color{white}\textbf{Segment} & \rowcolor{BleuNuit}\color{white}\textbf{Volume annuel (EVP)} & \rowcolor{BleuNuit}\color{white}\textbf{Parts de Marché cibles} & \rowcolor{BleuNuit}\color{white}\textbf{Nombre de dispositifs} \\
\midrule
\endfirsthead
\multicolumn{4}{c}{\tablename\ \thetable{} -- Suite de la page précédente} \\
\toprule
\rowcolor{BleuNuit}\color{white}\textbf{Segment} & \rowcolor{BleuNuit}\color{white}\textbf{Volume annuel (EVP)} & \rowcolor{BleuNuit}\color{white}\textbf{Parts de Marché cibles} & \rowcolor{BleuNuit}\color{white}\textbf{Nombre de dispositifs} \\
\midrule
\endhead
\midrule
\multicolumn{4}{r}{\textit{Suite à la page suivante...}} \\
\endfoot
\bottomrule
\multicolumn{4}{l}{\textit{Source : Les Promoteurs 2025-2026}} \\
\endlastfoot
Marché total Afrique centrale & 2 195 398 & — & — \\
\rowcolor{BleuClair}Cible ARGOS (Cameroun) & 755 398 & — & — \\
Année 1 (0,01\,\%) & — & 0,02\,\% & 200 \\
\rowcolor{BleuClair}Année 2 (0,02\,\%) & — & 0,06\,\% & 500 \\
Année 3 (0,04\,\%) & — & 0,13\,\% & 1000 \\
\end{xltabular}

\noindent Le Tableau~\ref{tab:marche_adressable} matérialise le potentiel de pénétration commerciale de la solution ARGOS face aux volumes massifs de conteneurs transitant par les corridors d'Afrique centrale et du Cameroun (plus de 750\,000 EVP identifiés sur le marché national cible). L'intérêt stratégique de ce tableau réside dans la progressivité réaliste de la croissance visée sur les trois premières années d'exploitation (passant d'une part de marché de 0,02\,\% avec 200 dispositifs actifs en Année 1, pour atteindre 0,13\,\% et 1\,000 unités en Année 3). Cette trajectoire graduée structure la feuille de route commerciale du projet, assurant un dimensionnement cohérent de la flotte connectée par rapport aux investissements et aux capacités logistiques à déployer.

\subsection{Modèle économique : le partage de la valeur ajoutée}
\label{sec:modele-economique}

Afin de lever les barrières à l’entrée et d’inciter les entreprises de transport maritime et d’import-export à adopter notre solution, nous avons retenu un modèle économique de type « Revenue Sharing » (partage des bénéfices) qui repose sur une logique gagnant-gagnant.

Pour illustrer le modèle de viabilité économique du projet ARGOS, nous avons défini un scénario opérationnel basé sur le transit de trois véhicules importés de Chine (Shijiazhuang) vers Douala (Cameroun).

\begin{enumerate}
    \item \textbf{Hypothèse opérationnelle :} Valeur de la cargaison : 50 000 000 FCFA.
    \item \textbf{Service de traçabilité (Transitaire) :} L’entreprise de transit facture au client une commission de 3\,\% de la valeur de la marchandise pour intégrer la solution de suivi ARGOS.
    \item \textbf{Répartition des revenus} (Commission de 3\,\% soit 1 500 000 FCFA) :
    \begin{itemize}
        \item \textbf{ARGOS (60\,\%) :} 900 000 FCFA, couvrant l'exploitation de la plateforme et le retour sur investissement.
        \item \textbf{Partenaires logistiques (40\,\%) :} 600 000 FCFA, destinés aux services opérationnels et à la maintenance du dispositif sur le terrain.
    \end{itemize}
\end{enumerate}

\begin{xltabular}{\textwidth}{@{} X r X X @{}}
\caption{Structure du modèle de partage} \label{tab:structure_partage} \\
\toprule
\rowcolor{BleuNuit}\color{white}\textbf{Éléments} & \rowcolor{BleuNuit}\color{white}\textbf{Montant} & \rowcolor{BleuNuit}\color{white}\textbf{Détail} & \rowcolor{BleuNuit}\color{white}\textbf{Justification} \\
\midrule
\endfirsthead
\multicolumn{4}{c}{\tablename\ \thetable{} -- Suite de la page précédente} \\
\toprule
\rowcolor{BleuNuit}\color{white}\textbf{Éléments} & \rowcolor{BleuNuit}\color{white}\textbf{Montant} & \rowcolor{BleuNuit}\color{white}\textbf{Détail} & \rowcolor{BleuNuit}\color{white}\textbf{Justification} \\
\midrule
\endhead
\midrule
\multicolumn{4}{r}{\textit{Suite à la page suivante...}} \\
\endfoot
\bottomrule
\multicolumn{4}{l}{\textit{Source : Les Promoteurs 2025-2026}} \\
\endlastfoot
Valeur du dispositif & 220 000 FCFA & Coût de production + marge & Basé sur coût unitaire 150 000 FCFA \\
\rowcolor{BleuClair}Facturation par trajet & 1 500 000 FCFA & 14,66\,\% & Modèle conservateur \\
Part ARGOS (60\,\%) & 900 000 FCFA & Amortissement, plateforme, marge & Financement du CAPEX \\
\rowcolor{BleuClair}Part entreprise (40\,\%) & 600 000 FCFA & — & — \\
- Maintenance (20\,\%) & 300 000 FCFA & Préventive et corrective & — \\
- Imprévus (20\,\%) & 300 000 FCFA & Sinistres, pertes, réparations & — \\
\end{xltabular}
Le tableau ci-dessus détaille le budget prévisionnel des frais d’établissement nécessaires à la création de l’entreprise.



% \\\\\\\\\\\\\\\\\\\\\\\\\\\\\\\\\\\\\\\\\\\\\\\\\\\\\\\\\\\\\\\\\\\\\\\\\\\\\\\\\\

\subsection{Inventaire des besoins}
La première étape de notre prévision financière consiste à inventorier l’ensemble des besoins nécessaires au lancement et à l’exploitation de notre activité. Ces besoins se décomposent en plusieurs postes d’investissement.

\subsubsection{Inventaire de frais de constitution}
Les frais de constitution regroupent l’ensemble des éléments qui composent les frais d’établissement nécessaires lors de la création ou l’acquisition de notre structure, afin de mettre en place son cadre juridique, fiscal et administratif indispensable à l’immatriculation de l’entreprise.

\begin{xltabular}{\textwidth}{@{} X r @{}}
\caption{Inventaire de frais de constitution} \label{tab:frais_constitution} \\
\toprule
\rowcolor{BleuNuit}\color{white}\textbf{Éléments} & \rowcolor{BleuNuit}\color{white}\textbf{Montant (FCFA)} \\
\midrule
\endfirsthead
\multicolumn{2}{c}{\tablename\ \thetable{} -- Suite de la page précédente} \\
\toprule
\rowcolor{BleuNuit}\color{white}\textbf{Éléments} & \rowcolor{BleuNuit}\color{white}\textbf{Montant (FCFA)} \\
\midrule
\endhead
\midrule
\multicolumn{2}{r}{\textit{Suite à la page suivante...}} \\
\endfoot
\bottomrule
\multicolumn{2}{l}{\textit{Source : Les Promoteurs 2025-2026}} \\
\endlastfoot
Immatriculation au RCCM & 41 500 \\
\rowcolor{BleuClair}Frais d'annonces legales & 23 000 \\
Frais notaries & 80 000 \\
Boite postale (creation) & 35 000 \\
\midrule
\textbf{TOTAL} & \textbf{179 500} \\
\end{xltabular}

Le tableau ci-dessus presente le modèle de partage des revenus par trajet, répartissant stratégiquement la facturation entre ARGOS (60,\%) pour l'amortissement technologique, et l'entreprise partenaire (40,\%) pour l'exploitation et la maintenance.

\subsubsection{Inventaire des besoins en immobilisations}
Les immobilisations sont des biens acquis par l’entreprise et qui sont utilisés de façon durable pour l’activité de l’entreprise. On distingue trois (3) catégories d’immobilisations : les immobilisations incorporelles, les immobilisations corporelles et les immobilisations financière

\paragraph{Les Immobilisations Incorporelles}
Elles représentent les actifs non monétaires qu’une entreprise utilise pour obtenir des avantages économiques futurs.
\begin{xltabular}{\textwidth}{@{} X c r @{}}
\caption{Inventaire de besoins en logiciels} \label{tab:besoins_logiciels} \\
\toprule
\rowcolor{BleuNuit}\color{white}\textbf{Eléments} & \rowcolor{BleuNuit}\color{white}\textbf{Quantité} & \rowcolor{BleuNuit}\color{white}\textbf{Montant (FCFA)} \\
\midrule
\endfirsthead
\multicolumn{3}{c}{\tablename\ \thetable{} -- Suite de la page précédente} \\
\toprule
\rowcolor{BleuNuit}\color{white}\textbf{Eléments} & \rowcolor{BleuNuit}\color{white}\textbf{Quantité} & \rowcolor{BleuNuit}\color{white}\textbf{Montant (FCFA)} \\
\midrule
\endhead
\midrule
\multicolumn{3}{r}{\textit{Suite à la page suivante...}} \\
\endfoot
\bottomrule
\multicolumn{3}{l}{\textit{Source : Les Promoteurs 2025-2026}} \\
\endlastfoot
Hebergement VPS & 1 & 180 000 \\
\rowcolor{BleuClair}Firebase Cloud Messaging & 1 & 120 000 \\
Certificat SSL & 1 & 25 000 \\
\rowcolor{BleuClair}Nom de Domaine & 1 & 15 000 \\
Windows Professionnels & 1 & 646 800 \\
API Maritime Visions & 1 & 275 000 \\
\midrule
\textbf{TOTAL} & & \textbf{1 261 800} \\
\end{xltabular}

Nous aurons besoin de certains logiciels pour le fonctionnement de notre activité tel que le tableau ci-dessus.

\paragraph{Les immobilisations corporelles}
Elles représentent les actifs physiques que l’entreprise entend utiliser sur plus d’un exercice comptable, elles sont destinées à être utilisées par l’entreprise dans le cadre de son activité, ou à être louées à des tiers ou encore être utilisées à des fins de gestion interne.
\begin{xltabular}{\textwidth}{@{} X c r @{}}
\caption{Agencement et Installations} \label{tab:agencement} \\
\toprule
\rowcolor{BleuNuit}\color{white}\textbf{Éléments} & \rowcolor{BleuNuit}\color{white}\textbf{Quantité} & \rowcolor{BleuNuit}\color{white}\textbf{Montant (FCFA)} \\
\midrule
\endfirsthead
\multicolumn{3}{c}{\tablename\ \thetable{} -- Suite de la page précédente} \\
\toprule
\rowcolor{BleuNuit}\color{white}\textbf{Éléments} & \rowcolor{BleuNuit}\color{white}\textbf{Quantité} & \rowcolor{BleuNuit}\color{white}\textbf{Montant (FCFA)} \\
\midrule
\endhead
\midrule
\multicolumn{3}{r}{\textit{Suite à la page suivante...}} \\
\endfoot
\bottomrule
\multicolumn{3}{l}{\textit{Source : Les Promoteurs 2025-2026}} \\
\endlastfoot
Agencement de Bureau & 1 & 800 000 \\
\rowcolor{BleuClair}Installation & 1 & 350 000 \\
Climatisations & 6 & 4 200 000 \\
\midrule
\textbf{Total} & & \textbf{5 350 000} \\
\end{xltabular}

Les travaux d’agencements serviront pour l’essentiel à rendre l’espace de travail fonctionnel, c’est-à-dire compatible avec la nature des activités du projet. On regroupe dans cette rubrique les travaux liés aux câblages électriques et informatiques, les décorations, les agencements des bureaux

\begin{xltabular}{\textwidth}{@{} X c r @{}}
\caption{Matériel et Outillages} \label{tab:materiel_outillages} \\
\toprule
\rowcolor{BleuNuit}\color{white}\textbf{Élément} & \rowcolor{BleuNuit}\color{white}\textbf{Quantité} & \rowcolor{BleuNuit}\color{white}\textbf{Montants (FCFA)} \\
\midrule
\endfirsthead
\multicolumn{3}{c}{\tablename\ \thetable{} -- Suite de la page précédente} \\
\toprule
\rowcolor{BleuNuit}\color{white}\textbf{Élément} & \rowcolor{BleuNuit}\color{white}\textbf{Quantité} & \rowcolor{BleuNuit}\color{white}\textbf{Montants (FCFA)} \\
\midrule
\endhead
\midrule
\multicolumn{3}{r}{\textit{Suite à la page suivante...}} \\
\endfoot
\bottomrule
\multicolumn{3}{l}{\textit{Source : Les Promoteurs 2025-2026}} \\
\endlastfoot
Onduleurs & 5 & 250 000 \\
\rowcolor{BleuClair}Regulateurs de tention & 6 & 392 400 \\
Groupe electrogene & 1 & 550 000 \\
\midrule
\textbf{TOTAL} & & \textbf{1 192 400} \\
\end{xltabular}

L’entreprise se dote d’équipements de protection et secours électriques susceptibles d’assurer la production de courant électrique en cas d’interruption de la fourniture d’énergie par le concessionnaire public, et ce afin de garder en état de fonctionnement permanent les équipements électriques

\begin{xltabular}{\textwidth}{@{} X c r @{}}
\caption{Matériel informatique} \label{tab:materiel_informatique} \\
\toprule
\rowcolor{BleuNuit}\color{white}\textbf{Éléments} & \rowcolor{BleuNuit}\color{white}\textbf{Quantité} & \rowcolor{BleuNuit}\color{white}\textbf{Montants (FCFA)} \\
\midrule
\endfirsthead
\multicolumn{3}{c}{\tablename\ \thetable{} -- Suite de la page précédente} \\
\toprule
\rowcolor{BleuNuit}\color{white}\textbf{Éléments} & \rowcolor{BleuNuit}\color{white}\textbf{Quantité} & \rowcolor{BleuNuit}\color{white}\textbf{Montants (FCFA)} \\
\midrule
\endhead
\midrule
\multicolumn{3}{r}{\textit{Suite à la page suivante...}} \\
\endfoot
\bottomrule
\multicolumn{3}{l}{\textit{Source : Les Promoteurs 2025-2026}} \\
\endlastfoot
Ordinateurs portables & 5 & 1 750 000 \\
\rowcolor{BleuClair}Star-link & 1 & 250 000 \\
Téléphone fixe & 6 & 120 000 \\
\rowcolor{BleuClair}Imprimante numérique & 3 & 225 000 \\
Rallonges & 6 & 30 000 \\
\rowcolor{BleuClair}Scanner & 1 & 100 000 \\
Câbles HDMI & 6 & 30 000 \\
Vidéo projecteur & 1 & 80 000 \\
\midrule
\textbf{TOTAL} & & \textbf{2 585 000} \\
\end{xltabular}

Il est prévu l’acquisition de matériel informatique et notamment de serveur pour un total de 2 585 000 FCFA. Trois configurations complètes d’ordinateurs seront acquises en An 1, et une en An 2.

\begin{xltabular}{\textwidth}{@{} X c r @{}}
\caption{Matériel de bureautique} \label{tab:materiel_bureautique} \\
\toprule
\rowcolor{BleuNuit}\color{white}\textbf{Éléments} & \rowcolor{BleuNuit}\color{white}\textbf{Quantité} & \rowcolor{BleuNuit}\color{white}\textbf{Montant (FCFA)} \\
\midrule
\endfirsthead
\multicolumn{3}{c}{\tablename\ \thetable{} -- Suite de la page précédente} \\
\toprule
\rowcolor{BleuNuit}\color{white}\textbf{Éléments} & \rowcolor{BleuNuit}\color{white}\textbf{Quantité} & \rowcolor{BleuNuit}\color{white}\textbf{Montant (FCFA)} \\
\midrule
\endhead
\midrule
\multicolumn{3}{r}{\textit{Suite à la page suivante...}} \\
\endfoot
\bottomrule
\multicolumn{3}{l}{\textit{Source : Les Promoteurs 2025-2026}} \\
\endlastfoot
Paquet de stylo & 5 & 25 000 \\
\rowcolor{BleuClair}Rame de format A4 & 10 & 200 000 \\
Chemises & 10 & 50 000 \\
\rowcolor{BleuClair}Agrafeuse & 5 & 25 000 \\
Paquet correcteur & 2 & 10 000 \\
\rowcolor{BleuClair}Classeur & 10 & 50 000 \\
Calculatrice & 3 & 60 000 \\
\midrule
\textbf{Total} & & \textbf{420 000} \\
\end{xltabular}

\begin{xltabular}{\textwidth}{@{} X c r @{}}
\caption{Matériel de bureau} \label{tab:materiel_bureau} \\
\toprule
\rowcolor{BleuNuit}\color{white}\textbf{Éléments} & \rowcolor{BleuNuit}\color{white}\textbf{Quantité} & \rowcolor{BleuNuit}\color{white}\textbf{Montants (FCFA)} \\
\midrule
\endfirsthead
\multicolumn{3}{c}{\tablename\ \thetable{} -- Suite de la page précédente} \\
\toprule
\rowcolor{BleuNuit}\color{white}\textbf{Éléments} & \rowcolor{BleuNuit}\color{white}\textbf{Quantité} & \rowcolor{BleuNuit}\color{white}\textbf{Montants (FCFA)} \\
\midrule
\endhead
\midrule
\multicolumn{3}{r}{\textit{Suite à la page suivante...}} \\
\endfoot
\bottomrule
\multicolumn{3}{l}{\textit{Source : Les Promoteurs 2025-2026}} \\
\endlastfoot
Tables de bureau & 5 & 300 000 \\
\rowcolor{BleuClair}Tables de réunions & 1 & 300 000 \\
Chaises de bureau & 5 & 200 000 \\
\rowcolor{BleuClair}Chaises de réunions & 10 & 250 000 \\
Chaises de visiteurs & 6 & 90 000 \\
\rowcolor{BleuClair}Extincteur & 3 & 75 000 \\
Armoire de bureau & 5 & 500 000 \\
\midrule
\textbf{Total} & & \textbf{1 715 000} \\
\end{xltabular}

Nous regroupons dans cette rubrique l’ensemble des immobilisations telles que meubles et autres équipements de bureau nécessaires à l’exploitation de l’entreprise. Un ensemble table et chaise de bureau sera acquis en An 2.

\begin{xltabular}{\textwidth}{@{} X c r @{}}
\caption{Récapitulatif des immobilisations} \label{tab:recap_immo} \\
\toprule
\rowcolor{BleuNuit}\color{white}\textbf{ELEMENT} & \rowcolor{BleuNuit}\color{white}\textbf{QUANTITÉ} & \rowcolor{BleuNuit}\color{white}\textbf{MONTANT (FCFA)} \\
\midrule
\endfirsthead
\multicolumn{3}{c}{\tablename\ \thetable{} -- Suite de la page précédente} \\
\toprule
\rowcolor{BleuNuit}\color{white}\textbf{ELEMENT} & \rowcolor{BleuNuit}\color{white}\textbf{QUANTITÉ} & \rowcolor{BleuNuit}\color{white}\textbf{MONTANT (FCFA)} \\
\midrule
\endhead
\midrule
\multicolumn{3}{r}{\textit{Suite à la page suivante...}} \\
\endfoot
\bottomrule
\multicolumn{3}{l}{\textit{Source : Les Promoteurs 2025-2026}} \\
\endlastfoot
Logiciel & 1 & 1 261 800 \\
\rowcolor{BleuClair}Tables de bureau & 5 & 300 000 \\
Tables de réunions & 1 & 300 000 \\
\rowcolor{BleuClair}Chaises de bureau & 5 & 200 000 \\
Chaises de réunions & 10 & 250 000 \\
\rowcolor{BleuClair}Chaises de visiteurs & 6 & 90 000 \\
Armoire de bureau & 5 & 500 000 \\
\rowcolor{BleuClair}Extinteurs & 3 & 75 000 \\
Ordinateurs portables & 5 & 1 750 000 \\
\rowcolor{BleuClair}Star-link & 1 & 250 000 \\
Téléphone fixe & 6 & 120 000 \\
\rowcolor{BleuClair}Imprimante numérique & 3 & 225 000 \\
Rallonges & 6 & 30 000 \\
\rowcolor{BleuClair}Scanner & 1 & 100 000 \\
Câbles HDMI & 6 & 30 000 \\
\rowcolor{BleuClair}Vidéo projecteur & 1 & 80 000 \\
Agencement de Bureau & 1 & 800 000 \\
\rowcolor{BleuClair}Installation & 1 & 350 000 \\
Climatisations & 6 & 4 200 000 \\
\rowcolor{BleuClair}Onduleurs & 5 & 250 000 \\
Regulateurs de tention & 6 & 392 400 \\
Groupe electrogene & 1 & 550 000 \\
\midrule
\textbf{Total immobilisation} & & \textbf{12 104 200} \\
\end{xltabular}

L'acquisition des immobilisations sera étalée sur les deux premières années. Le total des immobilisations à acquérir est de \textbf{12 104 200 FCFA} au cours de la première année d'investissement.

\subsection{Besoins en fonds de roulement (BFR)}
Le besoin en fonds de roulement est la somme d'argent dont une entreprise a besoin pour financer son cycle afin de couvrir les décalages de trésoreries entre les entrées et les sorties d'argent.
\begin{xltabular}{\textwidth}{@{} X r c r @{}}
\caption{Besoin en fonds de roulement (BFR) de démarrage (6 mois)} \label{tab:bfr_demarrage} \\
\toprule
\rowcolor{BleuNuit}\color{white}\textbf{Élément} & \rowcolor{BleuNuit}\color{white}\textbf{Charge Annuelle} & \rowcolor{BleuNuit}\color{white}\textbf{Mois} & \rowcolor{BleuNuit}\color{white}\textbf{Montant Préfinancé (FCFA)} \\
\midrule
\endfirsthead
\multicolumn{4}{c}{\tablename\ \thetable{} -- Suite de la page précédente} \\
\toprule
\rowcolor{BleuNuit}\color{white}\textbf{Élément} & \rowcolor{BleuNuit}\color{white}\textbf{Charge Annuelle} & \rowcolor{BleuNuit}\color{white}\textbf{Mois} & \rowcolor{BleuNuit}\color{white}\textbf{Montant Préfinancé (FCFA)} \\
\midrule
\endhead
\midrule
\multicolumn{4}{r}{\textit{Suite à la page suivante...}} \\
\endfoot
\bottomrule
\multicolumn{4}{l}{\textit{Source : Les Promoteurs 2025-2026}} \\
\endlastfoot
Stock de demarrage & 8 705 600 & 6 & 4 352 800 \\
\rowcolor{BleuClair}Abonnement fibre optique & 250 000 & 6 & 125 000 \\
Eau et Électricité & 1 080 000 & 6 & 540 000 \\
\rowcolor{BleuClair}Location de Bureau & 2 400 000 & 6 & 1 200 000 \\
Transport urbains & 150 000 & 6 & 75 000 \\
\rowcolor{BleuClair}Entretiens et reparations & 200 000 & 6 & 100 000 \\
Publicite & 600 000 & 6 & 300 000 \\
\rowcolor{BleuClair}Fournitures de bureau & 200 000 & 6 & 100 000 \\
Frais de communication & 50 000 & 6 & 25 000 \\
\rowcolor{BleuClair}Impot et taxes & 100 000 & 6 & 50 000 \\
Charge de personnel & 1 950 000 & 6 & 975 000 \\
\midrule
\textbf{Total des charges à préfinancer} & \textbf{15 685 600} & & \textbf{7 842 800} \\
\end{xltabular}

L'entreprise au moment de son démarrage doit engager certaines charges avant même de commencer la commercialisation de ses produits. Et même lorsqu'elle aura réalisé les premières services, on doute que ces premiers revenus soient suffisants pour soutenir toutes les charges d'exploitation générées. Aussi serait-il prudent de regrouper au titre de fonds de roulement de démarrage, les besoins en financement qui pourraient résulter du décalage entre les premières recettes et dépenses d'exploitation.
Nous retenons quelques charges annuelles et nous décidons de ne financer que la moitié d’entre elles au titre du besoin en fonds de roulement de démarrage : ce qui nous donne un montant de besoin en fonds de roulement de \textbf{7 842 800 FCFA}.

\begin{xltabular}{\textwidth}{@{} X r @{}}
\caption{Récapitulatif des besoins} \label{tab:recap_besoins} \\
\toprule
\rowcolor{BleuNuit}\color{white}\textbf{Élément} & \rowcolor{BleuNuit}\color{white}\textbf{Montant (FCFA)} \\
\midrule
\endfirsthead
\multicolumn{2}{c}{\tablename\ \thetable{} -- Suite de la page précédente} \\
\toprule
\rowcolor{BleuNuit}\color{white}\textbf{Élément} & \rowcolor{BleuNuit}\color{white}\textbf{Montant (FCFA)} \\
\midrule
\endhead
\midrule
\multicolumn{2}{r}{\textit{Suite à la page suivante...}} \\
\endfoot
\bottomrule
\multicolumn{2}{l}{\textit{Source : Les Promoteurs 2025-2026}} \\
\endlastfoot
Logiciel & 1 261 800 \\
\rowcolor{BleuClair}Amenagement et Installation & 5 350 000 \\
Materiel et Outillage & 1 192 400 \\
\rowcolor{BleuClair}Materiel informatique & 2 585 000 \\
Mobilier de bureau & 1 715 000 \\
\midrule
\textbf{Total immobilisation} & \textbf{12 104 200} \\
Besoin en fonds de roulement de demarrage & 7 842 800 \\
\midrule
\textbf{Total BFR de demarrage} & \textbf{19 947 000} \\
\end{xltabular}

\subsection{Plan d'amortissement des immobilisations}
Un plan d'amortissement est un tableau financier visant à présenter la dépréciation qui affectera une immobilisation du fait de l'usure et du temps. Afin de mesurer dans la durée l'amortissement, un plan est établi dès que le bien est porté à l'actif du bilan comptable. Ces immobilisations sont amorties conformément aux dispositions du code général des impôts du Cameroun à savoir les taux d'amortissement prévus et la méthode d'amortissement en vigueur c'est-à-dire l'amortissement linéaire.
Notre plan d'amortissement est établi sur une période de 3 ans. Le montant annuel des dotations reflète bien l'étalement dans le temps du programme d'acquisition des immobilisations tel qu'exprimé plus haut.

\begin{xltabular}{\textwidth}{@{} X r r c r r r @{}}
\caption{Amortissement des Immobilisations} \label{tab:amortissement_immo} \\
\toprule
\rowcolor{BleuNuit}\color{white}\textbf{Éléments} & \rowcolor{BleuNuit}\color{white}\textbf{Montants} & \rowcolor{BleuNuit}\color{white}\textbf{Taux} & \rowcolor{BleuNuit}\color{white}\textbf{Durée} & \rowcolor{BleuNuit}\color{white}\textbf{Année N} & \rowcolor{BleuNuit}\color{white}\textbf{Année N+1} & \rowcolor{BleuNuit}\color{white}\textbf{Année N+2} \\
\midrule
\endfirsthead
\multicolumn{7}{c}{\tablename\ \thetable{} -- Suite de la page précédente} \\
\toprule
\rowcolor{BleuNuit}\color{white}\textbf{Éléments} & \rowcolor{BleuNuit}\color{white}\textbf{Montants} & \rowcolor{BleuNuit}\color{white}\textbf{Taux} & \rowcolor{BleuNuit}\color{white}\textbf{Durée} & \rowcolor{BleuNuit}\color{white}\textbf{Année N} & \rowcolor{BleuNuit}\color{white}\textbf{Année N+1} & \rowcolor{BleuNuit}\color{white}\textbf{Année N+2} \\
\midrule
\endhead
\midrule
\multicolumn{7}{r}{\textit{Suite à la page suivante...}} \\
\endfoot
\bottomrule
\multicolumn{7}{l}{\textit{Source : Code général des impôts.}} \\
\endlastfoot
Logiciel & 1 261 800 & 33\,\% & 3 & 420 600 & 420 600 & 420 600 \\
\rowcolor{BleuClair}Agencement de Bureau & 800 000 & 20\,\% & 5 & 160 000 & 160 000 & 160 000 \\
Installation & 350 000 & 10\,\% & 10 & 35 000 & 35 000 & 35 000 \\
\rowcolor{BleuClair}Climatisations & 4 200 000 & 10\,\% & 10 & 420 000 & 420 000 & 420 000 \\
Tables de bureau & 300 000 & 10\,\% & 10 & 30 000 & 30 000 & 30 000 \\
\rowcolor{BleuClair}Tables de réunions & 300 000 & 10\,\% & 10 & 30 000 & 30 000 & 30 000 \\
Chaises de bureau & 200 000 & 10\,\% & 10 & 20 000 & 20 000 & 20 000 \\
\rowcolor{BleuClair}Chaises de réunions & 250 000 & 10\,\% & 10 & 25 000 & 25 000 & 25 000 \\
Chaises de visiteurs & 90 000 & 10\,\% & 10 & 9 000 & 9 000 & 9 000 \\
\rowcolor{BleuClair}Armoire de bureau & 500 000 & 10\,\% & 10 & 50 000 & 50 000 & 50 000 \\
Extincteur & 75 000 & 10\,\% & 10 & 7 500 & 7 500 & 7 500 \\
\rowcolor{BleuClair}Ordinateurs portables & 1 750 000 & 25\,\% & 4 & 437 500 & 437 500 & 437 500 \\
Star-link & 225 000 & 25\,\% & 4 & 56 250 & 56 250 & 56 250 \\
\rowcolor{BleuClair}Téléphone fixe & 120 000 & 25\,\% & 4 & 30 000 & 30 000 & 30 000 \\
Imprimante numérique & 225 000 & 25\,\% & 4 & 56 250 & 56 250 & 56 250 \\
\rowcolor{BleuClair}Rallonges & 30 000 & 10\,\% & 10 & 3 000 & 3 000 & 3 000 \\
Scanner & 100 000 & 25\,\% & 4 & 25 000 & 25 000 & 25 000 \\
\rowcolor{BleuClair}Câbles HDMI & 30 000 & 10\,\% & 10 & 3 000 & 3 000 & 3 000 \\
Vidéo projecteur & 80 000 & 25\,\% & 4 & 20 000 & 20 000 & 20 000 \\
\rowcolor{BleuClair}Onduleurs & 250 000 & 10\,\% & 10 & 25 000 & 25 000 & 25 000 \\
Regulateurs de tention & 392 400 & 10\,\% & 10 & 39 240 & 39 240 & 39 240 \\
Groupe electrogene & 550 000 & 10\,\% & 10 & 55 000 & 55 000 & 55 000 \\
\midrule
\textbf{TOTAL} & \textbf{12 079 200} & & \textbf{Total} & \textbf{1 957 340} & \textbf{1 957 340} & \textbf{1 957 340} \\
\end{xltabular}

Notre plan d’amortissement est établi sur une durée de 3 ans. Les immobilisations seront amorties pendant la première année, pour un total de 1 957 340 FCFA. 

\subsection{États financiers prévisionnels}
Un état financier est une représentation financière structurée des évènements affectant une entreprise et des différentes transactions effectuées par celle-ci. Son objectif est de fournir des informations sur la situation financière utile à la détermination du niveau de performance de l’entreprise et à l’évaluation de sa capacité d’autofinancement.
\subsection{Apport financier}
L'apport financier c'est l'ensemble des biens que les associés mettent à la disposition de la société en vue d'une exploitation commune. En contrepartie, les apporteurs reçoivent des titres soumis aux aléas de la société\cite{bpifrance}.
\begin{xltabular}{\textwidth}{@{} X X r @{}}
\caption{Source de financement} \label{tab:source_financement} \\
\toprule
\rowcolor{BleuNuit}\color{white}\textbf{Besoin} & \rowcolor{BleuNuit}\color{white}\textbf{Postes} & \rowcolor{BleuNuit}\color{white}\textbf{Montant (FCFA)} \\
\midrule
\endfirsthead
\multicolumn{3}{c}{\tablename\ \thetable{} -- Suite de la page précédente} \\
\toprule
\rowcolor{BleuNuit}\color{white}\textbf{Besoin} & \rowcolor{BleuNuit}\color{white}\textbf{Postes} & \rowcolor{BleuNuit}\color{white}\textbf{Montant (FCFA)} \\
\midrule
\endhead
\midrule
\multicolumn{3}{r}{\textit{Suite à la page suivante...}} \\
\endfoot
\bottomrule
\multicolumn{3}{l}{\textit{Source : Les promoteurs 2025-2026}} \\
\endlastfoot
MEKE NANILAOU B. M. & Gérant / Chef de projet & 4 500 000 \\
\rowcolor{BleuClair}KENNE METEU J. L. & Responsable technique & 4 100 000 \\
SIRRI NEH PRECIOUS D. N. & Resp. des affaires financières & 4 000 000 \\
\rowcolor{BleuClair}MAPPOU ZOUA A. D. & Responsable Marketing et Commerial & 3 700 000 \\
ATANGANA NJOH M. F. & Responsable Informatique & 3 700 000 \\
INVESTISSEUR & — & 10 000 000 \\
\midrule
\textbf{TOTAL} & & \textbf{30 000 000} \\
\end{xltabular}

Les membres du groupe vont apporter leur capital suivie d'un investisseur avec un solde de 10 000 000 FCFA.

\begin{xltabular}{\textwidth}{@{} X r r r r @{}}
\caption{Structuration de l’investissement} \label{tab:structuration_investissement} \\
\toprule
\rowcolor{BleuNuit}\color{white}\textbf{Designation} & \rowcolor{BleuNuit}\color{white}\textbf{Capital} & \rowcolor{BleuNuit}\color{white}\textbf{(\%)} & \rowcolor{BleuNuit}\color{white}\textbf{Emprunt} & \rowcolor{BleuNuit}\color{white}\textbf{Total} \\
\midrule
\endfirsthead
\multicolumn{5}{c}{\tablename\ \thetable{} -- Suite de la page précédente} \\
\toprule
\rowcolor{BleuNuit}\color{white}\textbf{Designation} & \rowcolor{BleuNuit}\color{white}\textbf{Capital} & \rowcolor{BleuNuit}\color{white}\textbf{(\%)} & \rowcolor{BleuNuit}\color{white}\textbf{Emprunt} & \rowcolor{BleuNuit}\color{white}\textbf{Total} \\
\midrule
\endhead
\midrule
\multicolumn{5}{r}{\textit{Suite à la page suivante...}} \\
\endfoot
\bottomrule
\multicolumn{5}{l}{\textit{Source : Les promoteurs 2025-2026}} \\
\endlastfoot
MEKE NANILAOU B. M. & 4 500 000 & 22,5\,\% & 0 & 4 500 000 \\
\rowcolor{BleuClair}KENNE METEU J. L. & 4 100 000 & 20,5\,\% & 0 & 4 100 000 \\
SIRRI NEH PRECIOUS D. N. & 4 000 000 & 20,0\,\% & 0 & 4 000 000 \\
\rowcolor{BleuClair}MAPPOU ZOUA A. D. & 3 700 000 & 18,5\,\% & 0 & 3 700 000 \\
ATANGANA NJOH M. F. & 3 700 000 & 18,5\,\% & 0 & 3 700 000 \\
INVESTISSEUR & 0 & 0\,\% & 10 000 000 & 10 000 000 \\
\midrule
\textbf{TOTAL} & \textbf{20 000 000} & \textbf{100\,\%} & \textbf{10 000 000} & \textbf{30 000 000} \\
\end{xltabular}

\subsection{Bilan d'ouverture}
Selon les normes IFRS (International Financial Reporting Standards) « Le bilan (état de la situation financière) présente les actifs, les passifs et les capitaux propres à une date donnée. Il reflète les ressources contrôlées et les obligations de l'entreprise. »
\begin{xltabular}{\textwidth}{@{} X r X r @{}}
\caption{Bilan d'ouverture de l’année 1} \label{tab:bilan_ouverture} \\
\toprule
\rowcolor{BleuNuit}\color{white}\textbf{ACTIF} & \rowcolor{BleuNuit}\color{white}\textbf{Montant} & \rowcolor{BleuNuit}\color{white}\textbf{PASSIF} & \rowcolor{BleuNuit}\color{white}\textbf{Montant} \\
\midrule
\endfirsthead
\multicolumn{4}{c}{\tablename\ \thetable{} -- Suite de la page précédente} \\
\toprule
\rowcolor{BleuNuit}\color{white}\textbf{ACTIF} & \rowcolor{BleuNuit}\color{white}\textbf{Montant} & \rowcolor{BleuNuit}\color{white}\textbf{PASSIF} & \rowcolor{BleuNuit}\color{white}\textbf{Montant} \\
\midrule
\endhead
\midrule
\multicolumn{4}{r}{\textit{Suite à la page suivante...}} \\
\endfoot
\bottomrule
\multicolumn{4}{l}{\textit{Source : Les promoteurs 2025-2026}} \\
\endlastfoot
Logiciels, services cloud & 1 261 800 & Capital Social & 20 000 000 \\
\rowcolor{BleuClair}\textbf{Immo. Incorporelles (A)} & \textbf{1 261 800} & \textbf{Capitaux propres} & \textbf{20 000 000} \\
Agencements et installation & 5 350 000 & Emprunt & 10 000 000 \\
\rowcolor{BleuClair}Matériel de bureau & 1 715 000 & \textbf{Dettes Financières} & \textbf{10 000 000} \\
Matériel bureautiques & 420 000 & & \\
\rowcolor{BleuClair}Matériel et Outillages & 1 192 400 & & \\
Matériel informatique & 2 585 000 & & \\
\rowcolor{BleuClair}\textbf{Immo. corporelles (B)} & \textbf{11 262 400} & & \\
Immobilisations financières & 0 & & \\
\midrule
\textbf{Actifs Immobilisés (A+B)} & \textbf{12 524 200} & \textbf{Ressources Stables} & \textbf{30 000 000} \\
\midrule
Stocks & 0 & Dettes Fournisseur & 0 \\
\rowcolor{BleuClair}Créances clients & 0 & Dettes sociales & 0 \\
\textbf{Actifs circulants} & \textbf{0} & \textbf{Passifs circulants} & \textbf{0} \\
\midrule
Caisse & 1 790 200 & & \\
\rowcolor{BleuClair}Banque & 15 685 600 & & \\
TOTAL ACTIFS & 30 000 000 & TOTAL PASSIFS & 30 000 000 \\
\end{xltabular}

Le bilan d'ouverture est équilibré en emplois et ressources à hauteur de 30 000 000 FCFA.
Les ressources sont constituées des apports en capital à hauteur de 20 000 000 FCFA et l'emprunt à mettre en place d'un montant de 10 000 000 FCFA. Ces ressources vont financer les immobilisations à acquérir en An 1 pour un montant de 12 104 200 FCFA. Ainsi que le montant du BFR (7 842 800 FCFA).

\section{Analyse financière de la viabilite du projet}

L'analyse financière est une étude qui consiste à tirer des renseignements sur la santé financière d'une entreprise, notamment en matière de solvabilité et de rentabilité, à partir des informations fournies par les états financiers. Faire une étude de la rentabilité de notre entreprise revient à faire un ensemble de calculs qui nous permettront de dire d'une part si le projet sera rentable et d'autre part, de voir si l'entreprise sera pérenne.

\subsubsection{Hypothèses}
Cette section énonce les estimations fondamentales concernant les coûts de production, les prix de service et les volumes de production.

\paragraph{Estimation des coûts et prix unitaires}
Nous détaillons ici les estimations des coûts de production et des prix unitaires d’un service des dispositifs et des abonnements comme suite : 

\begin{itemize}
    \item Estimation des coûts : 
    \begin{itemize} 
        \item Coût unitaire de production du dispositif à 220 000 FCFA. 
        \item Coût annuel d'abonnement 200 000 FCFA.
        \item Prix unitaire d’un service de 900 000 FCFA (cf. section~\ref{sec:modele-economique}) ;
        \item Prix du SAV de 50 000 FCFA.
    \end{itemize}
    \item Coût unitaire de production :
    
    « Le coût unitaire de production est le rapport entre le coût total engagé pour une production et le nombre d'unités produites. Il permet d'évaluer la performance économique d'un processus industriel. » (Comptabilité de gestion de Raymond Barre)
    
    \begin{xltabular}{\textwidth}{@{} X r @{}}
    \caption{Coût de production} \label{tab:cout_production} \\
    \toprule
    \rowcolor{BleuNuit}\color{white}\textbf{Éléments} & \rowcolor{BleuNuit}\color{white}\textbf{Montant (FCFA)} \\
    \midrule
    \endfirsthead
    \multicolumn{2}{c}{\tablename\ \thetable{} -- Suite de la page précédente} \\
    \toprule
    \rowcolor{BleuNuit}\color{white}\textbf{Éléments} & \rowcolor{BleuNuit}\color{white}\textbf{Montant (FCFA)} \\
    \midrule
    \endhead
    \midrule
    \multicolumn{2}{r}{\textit{Suite à la page suivante...}} \\
    \endfoot
    \bottomrule
    \multicolumn{2}{l}{\textit{Source : Les promoteurs 2025-2026}} \\
    \endlastfoot
    Coût de production du dispositif & 220 000 \\
    \rowcolor{BleuClair}Coût du SAV & 50 000 \\
    Coût de l'abonnement & 200 000 \\
    \end{xltabular}
    
    \item Estimation des autres charges d’exploitation : 
    Les charges d’exploitation prévisionnelles sont l’ensemble d’éléments qui participent directement ou indirectement à la modélisation du métier de l’entreprise. Elles comprennent les loyers, les salaires, les frais d’entretien et de maintenance observés sur trois années. 
\end{itemize}

Les autres charges d'exploitation sont estimées comme suit :

\begin{xltabular}{\textwidth}{@{} X r r r @{}}
\caption{Charges d’exploitation prévisionnelles} \label{tab:charges_exploitation} \\
\toprule
\rowcolor{BleuNuit}\color{white}\textbf{Éléments} & \rowcolor{BleuNuit}\color{white}\textbf{Année N} & \rowcolor{BleuNuit}\color{white}\textbf{Année N+1} & \rowcolor{BleuNuit}\color{white}\textbf{Année N+2} \\
\midrule
\endfirsthead
\multicolumn{4}{c}{\tablename\ \thetable{} -- Suite de la page précédente} \\
\toprule
\rowcolor{BleuNuit}\color{white}\textbf{Éléments} & \rowcolor{BleuNuit}\color{white}\textbf{Année N} & \rowcolor{BleuNuit}\color{white}\textbf{Année N+1} & \rowcolor{BleuNuit}\color{white}\textbf{Année N+2} \\
\midrule
\endhead
\midrule
\multicolumn{4}{r}{\textit{Suite à la page suivante...}} \\
\endfoot
\bottomrule
\multicolumn{4}{l}{\textit{Source : Les promoteurs 2025-2026}} \\
\endlastfoot
Abonnement fibre optique & 800 000 & 1 000 000 & 1 500 000 \\
\rowcolor{BleuClair}Mission et réception & 400 000 & 600 000 & 700 000 \\
Rémunération d'intermédiaires & 300 000 & 250 000 & 250 000 \\
\rowcolor{BleuClair}Eau \& Électricité & 720 000 & 720 000 & 720 000 \\
Imprévus & 50 000 & 50 000 & 50 000 \\
\rowcolor{BleuClair}Loyer & 1 740 000 & 1 740 000 & 1 740 000 \\
Publicité \& communication & 500 000 & 600 000 & 700 000 \\
\rowcolor{BleuClair}Transport urbains & 500 000 & 650 000 & 800 000 \\
Frais de Téléphone \& internet & 360 000 & 360 000 & 360 000 \\
\rowcolor{BleuClair}Frais de constitution & 179 500 & 0 & 0 \\
Impôts et taxes (patente, etc.) & 2 562 384 & 2 056 073 & 4 988 084 \\
\rowcolor{BleuClair}Fournitures de bureau & 420 000 & 420 000 & 420 000 \\
Entretien et réparations & 200 000 & 200 000 & 200 000 \\
\rowcolor{BleuClair}Charges de personnel & 18 010 200 & 19 464 870 & 20 088 300 \\
Services cloud & 500 000 & 600 000 & 750 000 \\
\rowcolor{BleuClair}Dotation aux Amortissement & 1 957 340 & 1 957 340 & 1 957 340 \\
Charge financière & 1 000 000 & 800 000 & 600 000 \\
\midrule
\textbf{Total des charges prévisionnelles} & \textbf{28 699 424} & \textbf{29 618 283} & \textbf{33 373 724} \\
\end{xltabular}

\subsection{Les prévisions financières}
Nous allons établir trois états financiers (compte de résultat, plan de financement et compte de bilan) qui nous permettront d'apprécier les performances de la société, sa situation financière, la qualité de son enrichissement.
\subsubsection{Formation du chiffre d'affaires prévisionnel}
\begin{xltabular}{\textwidth}{@{} X r r r r @{}}
\caption{Chiffres d’affaires annuels prévisionnels ARGOS} \label{tab:ca_previsionnel} \\
\toprule
\rowcolor{BleuNuit}\color{white}\textbf{Éléments} & \rowcolor{BleuNuit}\color{white}\textbf{Prix unitaire} & \rowcolor{BleuNuit}\color{white}\textbf{Année 1} & \rowcolor{BleuNuit}\color{white}\textbf{Année 2} & \rowcolor{BleuNuit}\color{white}\textbf{Année 3} \\
\midrule
\endfirsthead
\multicolumn{5}{c}{\tablename\ \thetable{} -- Suite de la page précédente} \\
\toprule
\rowcolor{BleuNuit}\color{white}\textbf{Éléments} & \rowcolor{BleuNuit}\color{white}\textbf{Prix unitaire} & \rowcolor{BleuNuit}\color{white}\textbf{Année 1} & \rowcolor{BleuNuit}\color{white}\textbf{Année 2} & \rowcolor{BleuNuit}\color{white}\textbf{Année 3} \\
\midrule
\endhead
\midrule
\multicolumn{5}{r}{\textit{Suite à la page suivante...}} \\
\endfoot
\bottomrule
\multicolumn{5}{l}{\textit{Source : Les promoteurs 2025-2026}} \\
\endlastfoot
Prix Unitaire du service & 900 000 & & & \\
\rowcolor{BleuClair}Nombre de boîtes & & 200 & 500 & 1 000 \\
\textbf{Chiffre d'affaire} & & \textbf{180 000 000} & \textbf{450 000 000} & \textbf{900 000 000} \\
\midrule
Prix Unitaire d'un trajet & 900 000 & & & \\
\rowcolor{BleuClair}Nombre de trajets & & 2 & 6 & 18 \\
\textbf{Chiffre d'affaire (trajet)} & & \textbf{1 800 000} & \textbf{5 400 000} & \textbf{16 200 000} \\
\midrule
Formation du personnel & 100 000 & & & \\
\rowcolor{BleuClair}Nombre de personnel & & 50 & 100 & 120 \\
\textbf{Chiffre d'affaire} & & \textbf{5 000 000} & \textbf{10 000 000} & \textbf{12 000 000} \\
\midrule
SAV & 50 000 & & & \\
\rowcolor{BleuClair}Nombre d'interventions & & 1 & 2 & 2 \\
\textbf{CHIFFRE D'AFFAIRE TOTAL} & & \textbf{186 850 000} & \textbf{465 500 000} & \textbf{928 300 000} \\
\end{xltabular}

\subsubsection{Formation des coûts de production prévisionnels}

\begin{xltabular}{\textwidth}{@{} X r r r @{}}
\caption{Formation des coûts de production} \label{tab:formation_couts} \\
\toprule
\rowcolor{BleuNuit}\color{white}\textbf{Éléments} & \rowcolor{BleuNuit}\color{white}\textbf{Année N} & \rowcolor{BleuNuit}\color{white}\textbf{Année N+1} & \rowcolor{BleuNuit}\color{white}\textbf{Année N+2} \\
\midrule
\endfirsthead
\multicolumn{4}{c}{\tablename\ \thetable{} -- Suite de la page précédente} \\
\toprule
\rowcolor{BleuNuit}\color{white}\textbf{Éléments} & \rowcolor{BleuNuit}\color{white}\textbf{Année N} & \rowcolor{BleuNuit}\color{white}\textbf{Année N+1} & \rowcolor{BleuNuit}\color{white}\textbf{Année N+2} \\
\midrule
\endhead
\midrule
\multicolumn{4}{r}{\textit{Suite à la page suivante...}} \\
\endfoot
\bottomrule
\multicolumn{4}{l}{\textit{Source : Les promoteurs 2025-2026}} \\
\endlastfoot
Coût de production & 44 000 000 & 110 000 000 & 220 000 000 \\
\rowcolor{BleuClair}Coût de l'abonnement & 40 000 000 & 100 000 000 & 200 000 000 \\
Quote part aux entreprises & 74 740 000 & 186 200 000 & 371 320 000 \\
\midrule
Coût de production total & 158 740 000 & 396 200 000 & 791 320 000 \\
\midrule
\textbf{Marge sur cout de production} & \textbf{28 110 000} & \textbf{69 300 000} & \textbf{136 980 000} \\
\end{xltabular}

Nous gardons un coût unitaire de production du dispositif constant sur l'ensemble de la période de prévision.

\subsubsection{Formation des charges de personnel}
(Burlaud, A., \& Simon, C. (2013). Comptabilité générale. 13e éd., Dunod) définissent les charges de personnel comme « l'ensemble des coûts liés à la rémunération du personnel salarié d'une entreprise, comprenant les salaires, les charges sociales et les avantages en nature ou en espèces. » Afin d'évaluer avec précision l'impact de ces dépenses sur la structure financière du projet ARGOS, la décomposition détaillée des différentes masses salariales s'articule selon la ventilation suivante :
\begin{itemize}
    \item {Masse salariale prévisionnelle}
L'organisation de la société se résume à 5 employés en année 1, nombre porté à 6 et 7 employés en années 2 et 3. Les promoteurs souhaitent promouvoir une politique de recrutement et de rémunération de ressources humaines de qualité. De même, les promoteurs pensent faire bénéficier au personnel une politique avantageuse (assurance maladie, mise en place d'une mutuelle des employés, …). Ils estiment à 10\,\% la totalité des avantages sociaux à accorder au personnel.
    \item Masse salariale brute : 
    Pour les premières années, la rémunération de l'équipe de direction a été fixée à un niveau modeste, notre priorité étant la viabilité et la croissance de l'entreprise qui n'est à ce stade qu'une start-up.
\end{itemize}

\begin{xltabular}{\textwidth}{@{} X c r r r r r r @{}}
\caption{Masse salariale prévisionnelle} \label{tab:masse_salariale} \\
\toprule
\rowcolor{BleuNuit}\color{white}\multirow{2}{*}{\textbf{Fonctions}} & \rowcolor{BleuNuit}\color{white}\multirow{2}{*}{\textbf{Eff.}} & \multicolumn{2}{c}{\rowcolor{BleuNuit}\color{white}\textbf{Année N}} & \multicolumn{2}{c}{\rowcolor{BleuNuit}\color{white}\textbf{Année N+1}} & \multicolumn{2}{c}{\rowcolor{BleuNuit}\color{white}\textbf{Année N+2}} \\
\cmidrule(lr){3-4} \cmidrule(lr){5-6} \cmidrule(lr){7-8}
& & \textbf{Mensuel} & \textbf{Annuel} & \textbf{Mensuel} & \textbf{Annuel} & \textbf{Mensuel} & \textbf{Annuel} \\
\midrule
\endfirsthead
\multicolumn{8}{c}{\tablename\ \thetable{} -- Suite de la page précédente} \\
\toprule
\rowcolor{BleuNuit}\color{white}\multirow{2}{*}{\textbf{Fonctions}} & \rowcolor{BleuNuit}\color{white}\multirow{2}{*}{\textbf{Eff.}} & \multicolumn{2}{c}{\rowcolor{BleuNuit}\color{white}\textbf{Année N}} & \multicolumn{2}{c}{\rowcolor{BleuNuit}\color{white}\textbf{Année N+1}} & \multicolumn{2}{c}{\rowcolor{BleuNuit}\color{white}\textbf{Année N+2}} \\
\cmidrule(lr){3-4} \cmidrule(lr){5-6} \cmidrule(lr){7-8}
& & \textbf{Mensuel} & \textbf{Annuel} & \textbf{Mensuel} & \textbf{Annuel} & \textbf{Mensuel} & \textbf{Annuel} \\
\midrule
\endhead
\midrule
\multicolumn{8}{r}{\textit{Suite à la page suivante...}} \\
\endfoot
\bottomrule
\multicolumn{8}{l}{\textit{Source : Les promoteurs 2025-2026}} \\
\endlastfoot
Gérant & 1 & 300 000 & 3 600 000 & 305 000 & 3 660 000 & 310 000 & 3 720 000 \\
\rowcolor{BleuClair}Resp. Technique & 1 & 250 000 & 3 000 000 & 275 000 & 3 300 000 & 285 000 & 3 420 000 \\
Responsable Financier & 1 & 250 000 & 3 000 000 & 275 000 & 3 300 000 & 285 000 & 3 420 000 \\
\rowcolor{BleuClair}Responsable Commercial & 1 & 250 000 & 3 000 000 & 275 000 & 3 300 000 & 285 000 & 3 420 000 \\
Responsable Informatique & 1 & 250 000 & 3 000 000 & 275 000 & 3 300 000 & 285 000 & 3 420 000 \\
\midrule
Total & 5 & 1 300 000 & 15 600 000 & 1 405 000 & 16 860 000 & 1 450 000 & 17 400 000 \\
\end{xltabular}

Comme l'indique le tableau, la masse salariale brute annuelle apres 3 ans est estimée à 17 400 000 FCFA, hors primes ou bonus éventuels.

\indent \paragraph{Retenues sociales et fiscales sur les salaires}
Conformément à la réglementation en vigueur, les différents prélèvements applicables sur les salaires bruts ont été rigoureusement intégrés dans nos prévisions financières. Notamment, la part salariale des cotisations sociales est calculée sur la base d'un taux de 4,2\,\% appliqué au salaire brut, portant le montant annuel cumulé de ces retenues pour l'ensemble du personnel à 730\,800 FCFA.

\begin{xltabular}{\textwidth}{@{} X r r r @{}}
\caption{Calcul des retenues CNPS} \label{tab:retenues_cnps} \\
\toprule
\rowcolor{BleuNuit}\color{white}\textbf{Postes} & \rowcolor{BleuNuit}\color{white}\textbf{Salaire brut (SBT)} & \rowcolor{BleuNuit}\color{white}\textbf{Retenus mensuelles} & \rowcolor{BleuNuit}\color{white}\textbf{Retenus annuelles} \\
\midrule
\endfirsthead
\multicolumn{4}{c}{\tablename\ \thetable{} -- Suite de la page précédente} \\
\toprule
\rowcolor{BleuNuit}\color{white}\textbf{Postes} & \rowcolor{BleuNuit}\color{white}\textbf{Salaire brut (SBT)} & \rowcolor{BleuNuit}\color{white}\textbf{Retenus mensuelles} & \rowcolor{BleuNuit}\color{white}\textbf{Retenus annuelles} \\
\midrule
\endhead
\midrule
\multicolumn{4}{r}{\textit{Suite à la page suivante...}} \\
\endfoot
\bottomrule
\multicolumn{4}{l}{\textit{Source : Les promoteurs 2025-2026}} \\
\endlastfoot
Gérant & 300 000 & 12 600 & 151 200 \\
\rowcolor{BleuClair}Responsable Technique & 250 000 & 10 500 & 126 000 \\
Responsable Financier & 250 000 & 10 500 & 126 000 \\
\rowcolor{BleuClair}Responsable Commercial & 250 000 & 10 500 & 126 000 \\
Responsable Informatique & 250 000 & 10 500 & 126 000 \\
\midrule
Total & 1 300 000 & 54 600 & 655 200 \\
\end{xltabular}

\indent Le calcul de l'Impôt sur le Revenu des Personnes Physiques (IRPP) s'articule en plusieurs étapes rigoureuses, appliquées conformément au barème en vigueur du Code Général des Impôts et majorées des Centimes Additionnels Communaux (CAC). Dans un premier temps, nous déterminons le Salaire Brut Taxable (SBT) sur lequel est appliqué un abattement forfaitaire annuel de 500\,000 FCFA avant l'application du taux dégressif de l'impôt. Ce mécanisme de calcul du Salaire Net Taxable (SNT) intègre ainsi la formule suivante : $\text{SNT} = (\text{SBT} - 30\,\% \times \text{SBT} - \text{CNPS})$.

\begin{xltabular}{\textwidth}{@{} X r r @{}}
\caption{Salaires nets taxables (SNT)} \label{tab:snt} \\
\toprule
\rowcolor{BleuNuit}\color{white}\textbf{Postes} & \rowcolor{BleuNuit}\color{white}\textbf{Mensuel (FCFA)} & \rowcolor{BleuNuit}\color{white}\textbf{Annuel (FCFA)} \\
\midrule
\endfirsthead
\multicolumn{3}{c}{\tablename\ \thetable{} -- Suite de la page précédente} \\
\toprule
\rowcolor{BleuNuit}\color{white}\textbf{Postes} & \rowcolor{BleuNuit}\color{white}\textbf{Mensuel (FCFA)} & \rowcolor{BleuNuit}\color{white}\textbf{Annuel (FCFA)} \\
\midrule
\endhead
\midrule
\multicolumn{3}{r}{\textit{Suite à la page suivante...}} \\
\endfoot
\bottomrule
\multicolumn{3}{l}{\textit{Source : Les promoteurs 2025-2026}} \\
\endlastfoot
Gérant & 197 400 & 2 368 800 \\
\rowcolor{BleuClair}Responsable technique & 164 500 & 1 974 000 \\
Responsable financier & 164 500 & 1 974 000 \\
\rowcolor{BleuClair}Responsable commercial & 164 500 & 1 974 000 \\
Responsable Informatique & 164 500 & 1 974 000 \\
\midrule
Total & 855 400 & 10 264 800 \\
\end{xltabular}

Le tableau suivant présente l'IRPP de base déterminé en se référant au barème du code général des impôts après l'abattement forfaitaire de 500 000 FCFA :

\begin{xltabular}{\textwidth}{@{} X r r r r r r @{}}
\caption{Détermination de l'IRPP de base plus le CAC} \label{tab:irpp} \\
\toprule
\rowcolor{BleuNuit}\color{white}\textbf{Postes} & \rowcolor{BleuNuit}\color{white}\textbf{SBT} & \rowcolor{BleuNuit}\color{white}\textbf{CNPS} & \rowcolor{BleuNuit}\color{white}\textbf{Abatt.} & \rowcolor{BleuNuit}\color{white}\textbf{SNT} & \rowcolor{BleuNuit}\color{white}\textbf{IRPP} & \rowcolor{BleuNuit}\color{white}\textbf{IRPP+CAC} \\
\midrule
\endfirsthead
\multicolumn{7}{c}{\tablename\ \thetable{} -- Suite de la page précédente} \\
\toprule
\rowcolor{BleuNuit}\color{white}\textbf{Postes} & \rowcolor{BleuNuit}\color{white}\textbf{SBT} & \rowcolor{BleuNuit}\color{white}\textbf{CNPS} & \rowcolor{BleuNuit}\color{white}\textbf{Abatt.} & \rowcolor{BleuNuit}\color{white}\textbf{SNT} & \rowcolor{BleuNuit}\color{white}\textbf{IRPP} & \rowcolor{BleuNuit}\color{white}\textbf{IRPP+CAC} \\
\midrule
\endhead
\midrule
\multicolumn{7}{r}{\textit{Suite à la page suivante...}} \\
\endfoot
\bottomrule
\multicolumn{7}{l}{\textit{Source : Les promoteurs 2025-2026}} \\
\endlastfoot
Gérant & 3 600 000 & 151 200 & 1 080 000 & 1 868 800 & 560 640 & 616 704 \\
\rowcolor{BleuClair}Resp. tech. & 3 000 000 & 126 000 & 900 000 & 1 474 000 & 442 200 & 486 420 \\
Resp. fin. & 3 000 000 & 126 000 & 900 000 & 1 474 000 & 442 200 & 486 420 \\
\rowcolor{BleuClair}Resp. com. & 3 000 000 & 126 000 & 900 000 & 1 474 000 & 442 200 & 486 420 \\
Resp. Inf. & 3 000 000 & 126 000 & 900 000 & 1 474 000 & 442 200 & 486 420 \\
\midrule
Total & 15 600 000 & 655 200 & 4 680 000 & 7 764 800 & 2 329 440 & 2 562 384 \\
\end{xltabular}

\indent Dans la suite de notre projet nous allons donc déterminer les salaires nets versés à l'ensemble du personnel en soustrayant les retenues CNPS et l'IRPP+CAC du salaire brut annuel. Le tableau ci-dessous présente le calcul détaillé des salaires nets pour chaque poste, ainsi que le total des salaires nets versés à l'ensemble du personnel.

\begin{xltabular}{\textwidth}{@{} X r r r r @{}}
\caption{Détermination des salaires nets} \label{tab:salaires_nets} \\
\toprule
\rowcolor{BleuNuit}\color{white}\textbf{Postes} & \rowcolor{BleuNuit}\color{white}\textbf{SBT (annuel)} & \rowcolor{BleuNuit}\color{white}\textbf{CNPS (4,2\,\%)} & \rowcolor{BleuNuit}\color{white}\textbf{IRPP 30\,\%} & \rowcolor{BleuNuit}\color{white}\textbf{Salaire net versé} \\
\midrule
\endfirsthead
\multicolumn{5}{c}{\tablename\ \thetable{} -- Suite de la page précédente} \\
\toprule
\rowcolor{BleuNuit}\color{white}\textbf{Postes} & \rowcolor{BleuNuit}\color{white}\textbf{SBT (annuel)} & \rowcolor{BleuNuit}\color{white}\textbf{CNPS (4,2\,\%)} & \rowcolor{BleuNuit}\color{white}\textbf{IRPP 30\,\%} & \rowcolor{BleuNuit}\color{white}\textbf{Salaire net versé} \\
\midrule
\endhead
\midrule
\multicolumn{5}{r}{\textit{Suite à la page suivante...}} \\
\endfoot
\bottomrule
\multicolumn{5}{l}{\textit{Source : Les promoteurs 2025-2026}} \\
\endlastfoot
Gérant & 3 600 000 & 151 200 & 616 704 & 2 832 096 \\
\rowcolor{BleuClair}Responsable Technique & 3 000 000 & 126 000 & 486 420 & 2 387 580 \\
Responsable financier & 3 000 000 & 126 000 & 486 420 & 2 387 580 \\
\rowcolor{BleuClair}Responsable Commercial & 3 000 000 & 126 000 & 486 420 & 2 387 580 \\
Responsable Informatique & 3 000 000 & 126 000 & 486 420 & 2 387 580 \\
\midrule
Total & 15 600 000 & 655 200 & 2 562 384 & 12 382 416 \\
\end{xltabular}

\begin{xltabular}{\textwidth}{@{} X r r r @{}}
\caption{Prévisions des charges sociales} \label{tab:charges_sociales} \\
\toprule
\rowcolor{BleuNuit}\color{white}\textbf{Éléments} & \rowcolor{BleuNuit}\color{white}\textbf{Année N} & \rowcolor{BleuNuit}\color{white}\textbf{Année N+1} & \rowcolor{BleuNuit}\color{white}\textbf{Année N+2} \\
\midrule
\endfirsthead
\multicolumn{4}{c}{\tablename\ \thetable{} -- Suite de la page précédente} \\
\toprule
\rowcolor{BleuNuit}\color{white}\textbf{Éléments} & \rowcolor{BleuNuit}\color{white}\textbf{Année N} & \rowcolor{BleuNuit}\color{white}\textbf{Année N+1} & \rowcolor{BleuNuit}\color{white}\textbf{Année N+2} \\
\midrule
\endhead
\midrule
\multicolumn{4}{r}{\textit{Suite à la page suivante...}} \\
\endfoot
\bottomrule
\multicolumn{4}{l}{\textit{Source : Les promoteurs 2025-2026}} \\
\endlastfoot
Salaire de base & 15 600 000 & 16 860 000 & 17 400 000 \\
\rowcolor{BleuClair}Allocation familiale (7\,\%) & 1 092 000 & 1 180 200 & 1 218 000 \\
Risque de travail (1,75\,\%) & 273 000 & 295 050 & 304 500 \\
\rowcolor{BleuClair}Accident de travail (2,5\,\%) & 390 000 & 421 500 & 435 000 \\
Pension viellesse (4,2\,\%) & 655 200 & 708 120 & 730 800 \\
\midrule
Total 1 & 2 410 200 & 2 604 870 & 2 688 300 \\
\midrule
\textbf{TOTAL Charge personnel} & \textbf{18 010 200} & \textbf{19 464 870} & \textbf{20 088 300} \\
\end{xltabular}

\subsection{Amortissement de l’emprunt et patente}

« L’amortissement d’un emprunt consiste à répartir le remboursement du capital emprunté sur une période donnée selon un échéancier défini. » (Colasse, B. (2011) Encyclopédie de la comptabilité, du contrôle de gestion et l’audit). 
Ainsi, nous avons prévu au terme de la première année de contracter un crédit auprès des investisseurs dont les caractéristiques se présentent comme suit : 

\begin{xltabular}{\textwidth}{@{} X r c r r r r @{}}
\caption{Amortissement de l’emprunt sur 5 ans} \label{tab:amortissement_emprunt} \\
\toprule
\rowcolor{BleuNuit}\color{white}\textbf{Année} & \rowcolor{BleuNuit}\color{white}\textbf{Dette début} & \rowcolor{BleuNuit}\color{white}\textbf{Taux} & \rowcolor{BleuNuit}\color{white}\textbf{Intérêt} & \rowcolor{BleuNuit}\color{white}\textbf{Amortis.} & \rowcolor{BleuNuit}\color{white}\textbf{Annuité} & \rowcolor{BleuNuit}\color{white}\textbf{Dette fin} \\
\midrule
\endfirsthead
\multicolumn{7}{c}{\tablename\ \thetable{} -- Suite de la page précédente} \\
\toprule
\rowcolor{BleuNuit}\color{white}\textbf{Année} & \rowcolor{BleuNuit}\color{white}\textbf{Dette début} & \rowcolor{BleuNuit}\color{white}\textbf{Taux} & \rowcolor{BleuNuit}\color{white}\textbf{Intérêt} & \rowcolor{BleuNuit}\color{white}\textbf{Amortis.} & \rowcolor{BleuNuit}\color{white}\textbf{Annuité} & \rowcolor{BleuNuit}\color{white}\textbf{Dette fin} \\
\midrule
\endhead
\midrule
\multicolumn{7}{r}{\textit{Suite à la page suivante...}} \\
\endfoot
\bottomrule
\multicolumn{7}{l}{\textit{Source : Les promoteurs 2025-2026}} \\
\endlastfoot
N & 10 000 000 & 10\,\% & 1 000 000 & 2 000 000 & 3 000 000 & 8 000 000 \\
\rowcolor{BleuClair}N+1 & 8 000 000 & 10\,\% & 800 000 & 2 000 000 & 2 800 000 & 6 000 000 \\
N+2 & 6 000 000 & 10\,\% & 600 000 & 2 000 000 & 2 600 000 & 4 000 000 \\
\rowcolor{BleuClair}N+3 & 4 000 000 & 10\,\% & 400 000 & 2 000 000 & 2 400 000 & 2 000 000 \\
N+4 & 2 000 000 & 10\,\% & 200 000 & 2 000 000 & 2 200 000 & 0 \\
\end{xltabular}
\textbf{Formules utilisées :}
\begin{itemize}
    \item Amortissement = Valeur de l’emprunt / 5 (ans)
    \item Intérêt = Valeur en début de période $\times$ Taux d’intérêt ;
    \item Mensualité = Amortissement + Intérêt ;
    \item Dette en fin de période = Valeur début de période – Amortissement.
\end{itemize}

\subsubsection{Patente}
La patente est un impôt direct dû par toute personne morale ou physique, du fait de l’exercice d’une activité économique, commerciale, industrielle ou toute profession à but lucratif. Cet impôt est perçu au profit des collectivités locales (mairies et FEICOM) et de la CRTV. Au Cameroun, elle est un impôt annuel payé au plus tard le 28 février de chaque année.

\begin{xltabular}{\textwidth}{@{} X l r r @{}}
\caption{Prévision de la patente à payer sur 3 ans} \label{tab:patente} \\
\toprule
\rowcolor{BleuNuit}\color{white}\textbf{Élément} & \rowcolor{BleuNuit}\color{white}\textbf{Calcul} & \rowcolor{BleuNuit}\color{white}\textbf{Année N+1} & \rowcolor{BleuNuit}\color{white}\textbf{Année N+2} \\
\midrule
\endfirsthead
\multicolumn{4}{c}{\tablename\ \thetable{} -- Suite de la page précédente} \\
\toprule
\rowcolor{BleuNuit}\color{white}\textbf{Élément} & \rowcolor{BleuNuit}\color{white}\textbf{Calcul} & \rowcolor{BleuNuit}\color{white}\textbf{Année N+1} & \rowcolor{BleuNuit}\color{white}\textbf{Année N+2} \\
\midrule
\endhead
\midrule
\multicolumn{4}{r}{\textit{Suite à la page suivante...}} \\
\endfoot
\bottomrule
\multicolumn{4}{l}{\textit{Source : Les promoteurs 2025-2026}} \\
\endlastfoot
Patente principal & CA (N-1) x 0,494\,\% & 923 039 & 2 299 570 \\
\rowcolor{BleuClair}Redevance audio. & 100\,\% du principal & 923 039 & 2 299 570 \\
FEICOM & 10\,\% du principal & 92 303,90 & 229 957 \\
\rowcolor{BleuClair}CAC & 3\,\% du principal & 27 691,17 & 68 987,10 \\
TDL & Barème & 90 000 & 90 000 \\
\midrule
\textbf{Total Patente} & & \textbf{2 056 073} & \textbf{4 988 084} \\
\end{xltabular}

\subsection{Compte d'exploitation prévisionnel}
« Le compte de résultat prévisionnel est un état financier présenté sous forme de liste, qui permet d'apprécier le niveau d'activité économique d'une entité sur une période donnée, généralement appelée exercice comptable. Il regroupe l'ensemble des produits et des charges généralement liés à cette période, et permet d'aboutir à un résultat net. »

\begin{xltabular}{\textwidth}{@{} X r r r @{}}
\caption{Compte de résultat prévisionnel} \label{tab:compte_resultat} \\
\toprule
\rowcolor{BleuNuit}\color{white}\textbf{Rubriques} & \rowcolor{BleuNuit}\color{white}\textbf{Année N} & \rowcolor{BleuNuit}\color{white}\textbf{Année N+1} & \rowcolor{BleuNuit}\color{white}\textbf{Année N+2} \\
\midrule
\endfirsthead
\multicolumn{4}{c}{\tablename\ \thetable{} -- Suite de la page précédente} \\
\toprule
\rowcolor{BleuNuit}\color{white}\textbf{Rubriques} & \rowcolor{BleuNuit}\color{white}\textbf{Année N} & \rowcolor{BleuNuit}\color{white}\textbf{Année N+1} & \rowcolor{BleuNuit}\color{white}\textbf{Année N+2} \\
\midrule
\endhead
\midrule
\multicolumn{4}{r}{\textit{Suite à la page suivante...}} \\
\endfoot
\bottomrule
\multicolumn{4}{l}{\textit{Source : Les promoteurs 2025-2026}} \\
\endlastfoot
Chiffre d’affaires net ARGOS & 186 850 000 & 465 500 000 & 928 300 000 \\
\rowcolor{BleuClair}Coût de production & 158 740 000 & 396 200 000 & 791 320 000 \\
Marge sur cout de production & 28 110 000 & 69 300 000 & 136 980 000 \\
\midrule
Abonnement fibre optique & 800 000 & 1 000 000 & 1 500 000 \\
\rowcolor{BleuClair}Eau \& électricité & 720 000 & 720 000 & 720 000 \\
Imprévus & 50 000 & 50 000 & 50 000 \\
\rowcolor{BleuClair}Loyer & 1 740 000 & 1 740 000 & 1 740 000 \\
Mission et réception & 400 000 & 600 000 & 700 000 \\
\rowcolor{BleuClair}Publicité \& communication & 500 000 & 600 000 & 700 000 \\
Rémunération d'intermédiaires & 300 000 & 250 000 & 250 000 \\
\rowcolor{BleuClair}Frais de téléphone \& internet & 360 000 & 360 000 & 360 000 \\
Frais de constitution & 179 500 & 0 & 0 \\
\rowcolor{BleuClair}Impôts et taxes (patente, etc.) & 2 562 384 & 2 056 073 & 4 988 084 \\
Transport urbains & 500 000 & 650 000 & 800 000 \\
\rowcolor{BleuClair}Fournitures de bureau & 420 000 & 420 000 & 420 000 \\
Entretien et réparations & 200 000 & 200 000 & 200 000 \\
\rowcolor{BleuClair}Services cloud & 500 000 & 600 000 & 750 000 \\
\midrule
Valeur ajoutée (VA) & 18 878 116 & 60 053 927 & 123 801 916 \\
\rowcolor{BleuClair}Charges de personnel & 18 010 200 & 19 464 870 & 20 088 300 \\
Excédent brut d’exploitation (EBE) & 867 916 & 40 589 057 & 103 713 616 \\
\midrule
\rowcolor{BleuClair}Dotation aux amortissements & 1 957 340 & 1 957 340 & 1 957 340 \\
Résultat d’exploitation & -1 089 424 & 38 631 717 & 101 756 276 \\
\rowcolor{BleuClair}Produits financiers & 0 & 0 & 0 \\
Charges financières & 1 000 000 & 800 000 & 600 000 \\
\rowcolor{BleuClair}Résultat financier & -1 000 000 & -800 000 & -600 000 \\
Résultat AO & -2 089 424 & 37 831 717 & 101 156 276 \\
\midrule
\rowcolor{BleuClair}Résultat courant avant impôts & -2 089 424 & 37 831 717 & 101 156 276 \\
Impôts sur les sociétés (27,5\,\%) & 0 & 10 403 722 & 27 817 976 \\
\midrule
\textbf{Résultat net} & \textbf{-2 089 424} & \textbf{27 427 995} & \textbf{73 338 300} \\
\end{xltabular}

\textbf{Formules utilisées :}
\begin{itemize}
    \item Valeur ajoutée (VA) = Chiffre d’affaires - total des charges
    \item Excédent brut d’exploitation (EBE) = VA - charges du personnel
    \item Résultat financier = produits financiers - charges financières
    \item Résultat des activités ordinaires (RAO) = résultat d’exploitation + résultat financier
    \item Résultat net = résultat courant avant impôts - impôts sur le résultat
\end{itemize}

\subsection{Variation du BFR et plan de financement prévisionnel}
« Les variations du besoin en fonds de roulement (BFR) prévisionnel traduisent les évolutions des besoins de financement liés au cycle d'exploitation de l'entreprise, notamment les délais entre les encaissements des ventes et les décaissements des achats. Ce poste financier prévisionnel permet d'anticiper les décalages de trésorerie dus aux créances clients, dettes fournisseurs et aux stocks. Une hausse du BFR signifie que l'entreprise immobilise davantage de ressources dans son exploitation, tandis qu'une baisse indique une amélioration de la liquidité. »\cite{vernimmen_dalloz,vernimmen_fnac}

\begin{xltabular}{\textwidth}{@{} X r r r @{}}
\caption{Variation en fonds de roulement} \label{tab:variation_bfr} \\
\toprule
\rowcolor{BleuNuit}\color{white}\textbf{Éléments de variation du BFR} & \rowcolor{BleuNuit}\color{white}\textbf{Année N} & \rowcolor{BleuNuit}\color{white}\textbf{Année N+1} & \rowcolor{BleuNuit}\color{white}\textbf{Année N+2} \\
\midrule
\endfirsthead
\multicolumn{4}{c}{\tablename\ \thetable{} -- Suite de la page précédente} \\
\toprule
\rowcolor{BleuNuit}\color{white}\textbf{Éléments de variation du BFR} & \rowcolor{BleuNuit}\color{white}\textbf{Année N} & \rowcolor{BleuNuit}\color{white}\textbf{Année N+1} & \rowcolor{BleuNuit}\color{white}\textbf{Année N+2} \\
\midrule
\endhead
\midrule
\multicolumn{4}{r}{\textit{Suite à la page suivante...}} \\
\endfoot
\bottomrule
\multicolumn{4}{l}{\textit{Source : Les promoteurs 2025-2026}} \\
\endlastfoot
Stocks (en jours) & 8 & 7 & 5 \\
\rowcolor{BleuClair}Creance clients (en jours) & 7 & 11 & 10 \\
Fournisseurs (en jours) & 3 & 5 & 7 \\
\rowcolor{BleuClair}\textbf{BFR (en jour)} & \textbf{12} & \textbf{13} & \textbf{8} \\
\midrule
Stocks (valeur) & 1 665 745 & 10 568 473 & 25 089 560 \\
\rowcolor{BleuClair}Creance clients (valeur) & 1 273 297 & 7 156 335 & 20 385 338 \\
Fournisseurs (valeur) & 548 806 & 695 236 & 686 008 \\
\midrule
\textbf{BFR (en valeur)} & \textbf{2 390 236} & \textbf{17 029 572} & \textbf{44 788 890} \\
\end{xltabular}

\subsubsection{Plan de financement prévisionnel}
Le plan de financement prévisionnel est un tableau financier structuré qui présente l'ensemble des ressources (apports, emprunts, subventions, autofinancement…) et des emplois (investissements, charges de démarrage, remboursement de dettes, BFR…) nécessaires au lancement et à la continuité d'un projet ou d'une activité sur une période donnée, en garantissant son équilibre financier tout au long de l'exercice comptable. Il constitue un outil essentiel d’aide à la décision pour les investisseurs et partenaires financiers.

\begin{xltabular}{\textwidth}{@{} X r r r @{}}
\caption{Plan de financement} \label{tab:plan_financement} \\
\toprule
\rowcolor{BleuNuit}\color{white}\textbf{Ressources} & \rowcolor{BleuNuit}\color{white}\textbf{Année N} & \rowcolor{BleuNuit}\color{white}\textbf{Année N+1} & \rowcolor{BleuNuit}\color{white}\textbf{Année N+2} \\
\midrule
\endfirsthead
\multicolumn{4}{c}{\tablename\ \thetable{} -- Suite de la page précédente} \\
\toprule
\rowcolor{BleuNuit}\color{white}\textbf{Ressources} & \rowcolor{BleuNuit}\color{white}\textbf{Année N} & \rowcolor{BleuNuit}\color{white}\textbf{Année N+1} & \rowcolor{BleuNuit}\color{white}\textbf{Année N+2} \\
\midrule
\endhead
\midrule
\multicolumn{4}{r}{\textit{Suite à la page suivante...}} \\
\endfoot
\bottomrule
\multicolumn{4}{l}{\textit{Source : Les promoteurs 2025-2026}} \\
\endlastfoot
Capital social & 20 000 000 & 20 000 000 & 20 000 000 \\
\rowcolor{BleuClair}Prime d'emissions et reserve & 0 & 0 & 0 \\
Emprunt & 10 000 000 & 0 & 0 \\
\rowcolor{BleuClair}Compte courant associe & 0 & 0 & 0 \\
Cashflow & 3 825 256 & 32 502 675 & 78 122 980 \\
\midrule
Total ressources & 33 825 256 & 52 502 675 & 98 122 980 \\
\midrule
Emplois & & & \\
\rowcolor{BleuClair}Investissements anciens & 0 & 0 & 0 \\
Investissements nouveaux & 12 079 200 & 0 & 0 \\
\rowcolor{BleuClair}Remboursement de l'emprunt & 8 000 000 & 6 000 000 & 4 000 000 \\
Remboursement de compte courant & 0 & 0 & 0 \\
\rowcolor{BleuClair}Dividendes distribues & 0 & 0 & 0 \\
Variation du BFR & 2 390 236 & 17 029 572 & 44 788 890 \\
\midrule
Total emplois & 22 469 436 & 23 029 572 & 48 788 890 \\
\midrule
Trésorerie de l'exercice & 11 355 820 & 29 473 103 & 49 334 090 \\
Trésorerie cumulée & 11 355 820 & 40 828 923 & 90 163 013 \\
\end{xltabular}

Pour boucler le programme d'investissement et couvrir les charges d'exploitation de la première année, il n'est pas prévu de financement bancaire, mais plutôt des apports en capital pour un montant de 20 000 000 FCFA, complété par un apport de fonds externes sous forme d’emprunt pour un montant de 10 000 000 FCFA. 

\subsection{Compte de bilan prévisionnel}
« Le compte de bilan est une photographie de l'entreprise à une date donnée, qui permet d'évaluer ses ressources, ses emplois et sa structure financière. » (Comprendre et interpréter un bilan comptable, Paul-Jacques Lehmann)

\begin{xltabular}{\textwidth}{@{} X r r r X r r r @{}}
\caption{Bilan prévisionnel} \label{tab:bilan_previsionnel} \\
\toprule
\rowcolor{BleuNuit}\color{white}\textbf{ACTIF} & \rowcolor{BleuNuit}\color{white}\textbf{N} & \rowcolor{BleuNuit}\color{white}\textbf{N+1} & \rowcolor{BleuNuit}\color{white}\textbf{N+2} & \rowcolor{BleuNuit}\color{white}\textbf{PASSIF} & \rowcolor{BleuNuit}\color{white}\textbf{N} & \rowcolor{BleuNuit}\color{white}\textbf{N+1} & \rowcolor{BleuNuit}\color{white}\textbf{N+2} \\
\midrule
\endfirsthead
\multicolumn{8}{c}{\tablename\ \thetable{} -- Suite de la page précédente} \\
\toprule
\rowcolor{BleuNuit}\color{white}\textbf{ACTIF} & \rowcolor{BleuNuit}\color{white}\textbf{N} & \rowcolor{BleuNuit}\color{white}\textbf{N+1} & \rowcolor{BleuNuit}\color{white}\textbf{N+2} & \rowcolor{BleuNuit}\color{white}\textbf{PASSIF} & \rowcolor{BleuNuit}\color{white}\textbf{N} & \rowcolor{BleuNuit}\color{white}\textbf{N+1} & \rowcolor{BleuNuit}\color{white}\textbf{N+2} \\
\midrule
\endhead
\midrule
\multicolumn{8}{r}{\textit{Suite à la page suivante...}} \\
\endfoot
\bottomrule
\multicolumn{8}{l}{\textit{Source : Les promoteurs 2025-2026}} \\
\endlastfoot
Logiciel & 841 200 & 420 600 & 0 & Capital social & 20 000 000 & 20 000 000 & 20 000 000 \\
\rowcolor{BleuClair}\textbf{Immo. Incorporelles} & \textbf{841 200} & \textbf{420 600} & \textbf{0} & Résultat net & -2 089 424 & 27 427 995 & 73 338 300 \\
Agencement \& Install. & 4 735 000 & 4 120 000 & 3 505 000 & RAN & 0 & -2 089 424 & 27 427 995 \\
\rowcolor{BleuClair}Matériel de bureau & 1 543 500 & 1 372 000 & 1 200 500 & \textbf{Capitaux propres} & \textbf{17 910 576} & \textbf{45 338 571} & \textbf{120 766 295} \\
Materiel \& Outillages & 1 073 160 & 953 920 & 834 680 & Emprunt & 8 000 000 & 6 000 000 & 4 000 000 \\
\rowcolor{BleuClair}Matériel informatique & 1 954 000 & 1 315 500 & 684 500 & \textbf{Dette financière} & \textbf{8 000 000} & \textbf{6 000 000} & \textbf{4 000 000} \\
\textbf{Immo. corporelles} & 9 305 660 & 7 761 420 & 6 224 680 & \textbf{Ressources stables} & \textbf{25 910 576} & \textbf{51 338 571} & \textbf{124 766 295} \\
\midrule
\rowcolor{BleuClair}\textbf{Actif Immobilise} & \textbf{10 146 860} & \textbf{8 182 020} & \textbf{6 224 680} & Fournisseurs & 548 806 & 695 236 & 686 008 \\
\midrule
Stocks & 1 665 745 & 10 568 473 & 25 089 560 & Dettes fiscales & 2 562 384 & 2 056 073 & 4 988 084 \\
\rowcolor{BleuClair}Créances Clients & 1 273 297 & 7 156 335 & 20 385 338 & \textbf{Passif circulant} & \textbf{3 111 190} & \textbf{2 751 309} & \textbf{5 674 092} \\
\textbf{Actifs circulants} & \textbf{2 939 042} & \textbf{17 724 808} & \textbf{45 474 898} & & & & \\
\midrule
Banque & 11 355 820 & 40 828 923 & 90 163 013 & & & & \\
\rowcolor{BleuClair}Caisse & 343 004 & 356 492 & 494 400 & & & & \\
TOTAL GENERAL & 24 784 726 & 67 092 243 & 142 356 991 & TOTAL GENERAL & 24 784 726 & 67 092 243 & 142 356 991 \\
\end{xltabular}

En l'état, les ressources propres du bilan sont pour l'essentiel constituées des montants contenus dans les postes Capital social, Report à nouveau, et Résultat. Elles augmentent continuellement jusqu'en An 3 pour atteindre le montant de 142,3 millions FCFA. Le bilan est équilibré en ressources et en emplois tout au long de la période de prévisions.

\subsection{Seuil de rentabilité et indicateurs de performance}
Le seuil de rentabilité constitue un outil de pilotage financier fondamental pour la pérennité du projet ARGOS. 
Il permet de déterminer le niveau d'activité minimal et exprimé en volume de chiffre d'affaires ou en nombre de prestations ce que l'entreprise doit réaliser pour couvrir l'intégralité de ses charges fixes et variables, 
marquant ainsi le point d'équilibre où le résultat net devient nul. 
L'analyse de cet indicateur clé, détaillée ci-après, offre aux promoteurs une visibilité indispensable sur la sécurité financière, 
le risque d'exploitation et la capacité de l'entreprise à générer des profits à court et moyen terme.

\subsubsection{Seuil de rentabilité et point mort}
Encore appelé chiffre d'affaires critique, le Seuil de Rentabilité (SR) est la valeur du chiffre d'affaires pour laquelle le résultat de l'entreprise est nul. Autrement dit, c'est le niveau du chiffre d'affaires à partir duquel l'entreprise ne réalise ni perte ni profit : elle peut juste couvrir l'ensemble de ses charges. Sa détermination passe par l'établissement d'un tableau différentiel d'exploitation (Guillaume Delemarle sur L-Expert-Comptable.com).

\textbf{Point mort :}
Le point mort (PM) quant à lui, représente la date exacte à partir de laquelle l'activité de l'entreprise commence à générer des profits. En intégrant des paramètres clés tels que le chiffre d'affaires critique (le seuil de rentabilité) et le volume des charges fixes, il permet de mesurer l'horizon temporel de la rentabilité. 

Les relations mathématiques fondamentales régissant ces indicateurs s'établissent selon les formules suivantes :
\begin{itemize}
    \item Seuil de rentabilité (SR) = Charges fixes (CF) / Taux de la marge sur coûts variables (TMCV)
    \item Taux de marge sur coûts variables (TMCV) = (CA – CV) / CA
    \item Point mort (en jours) = (SR / CA) $\times$ 360
\end{itemize}

Afin d'évaluer concrètement l'atteinte de cet équilibre financier pour le projet ARGOS, la déclinaison des calculs et des résultats s'articule selon le tableau ci-après :

\begin{xltabular}{\textwidth}{@{} X r r r @{}}
\caption{Seuil de rentabilité et point mort} \label{tab:seuil_rentabilite} \\
\toprule
\rowcolor{BleuNuit}\color{white}\textbf{Élément} & \rowcolor{BleuNuit}\color{white}\textbf{Année N} & \rowcolor{BleuNuit}\color{white}\textbf{Année N+1} & \rowcolor{BleuNuit}\color{white}\textbf{Année N+2} \\
\midrule
\endfirsthead
\multicolumn{4}{c}{\tablename\ \thetable{} -- Suite de la page précédente} \\
\toprule
\rowcolor{BleuNuit}\color{white}\textbf{Élément} & \rowcolor{BleuNuit}\color{white}\textbf{Année N} & \rowcolor{BleuNuit}\color{white}\textbf{Année N+1} & \rowcolor{BleuNuit}\color{white}\textbf{Année N+2} \\
\midrule
\endhead
\midrule
\multicolumn{4}{r}{\textit{Suite à la page suivante...}} \\
\endfoot
\bottomrule
\multicolumn{4}{l}{\textit{Source : Les promoteurs 2025-2026}} \\
\endlastfoot
Chiffre d’affaires net ARGOS & 186 850 000 & 465 500 000 & 928 300 000 \\
\rowcolor{BleuClair}Charges variables & 162 002 384 & 399 106 073 & 797 258 084 \\
Marge sur coûts variables & 24 847 616 & 66 393 927 & 131 041 916 \\
\rowcolor{BleuClair}Taux de marge & 0,133 & 0,143 & 0,141 \\
Charges fixes (CF) & 26 937 040 & 17 718 488 & 1 737 664 \\
\midrule
\textbf{Seuil de rentabilité (SR)} & \textbf{202 562 126} & \textbf{124 227 568} & \textbf{12 309 600} \\
\rowcolor{BleuClair}\textbf{Point mort (PM) en jours} & \textbf{390,27} & \textbf{96,07} & \textbf{4,77} \\
\end{xltabular}

Nous constatons que le chiffre d’affaires minimum à réaliser par notre entreprise pour ne pas perdre en termes d’argent ou être rentable en première année est de 202,5 millions FCFA, en deuxième année de 124,2 millions FCFA, et pour la troisième année de 123 millions FCFA. 
Le point mort (PM) est un indicateur complémentaire calculé à partir du seuil de rentabilité, permettant de savoir à quel moment l’entreprise devient rentable. D’après le tableau qui précède, notre entreprise atteindra son PM en année 1 dans 390 jours, en année 2 dans 96 jours, et en année 3 en 5 jours. En définitif, nous constatons que l’entreprise sera rentable dès le premier mois suivant son démarrage effectif.

\subsection{Cash-flow et VAN}
Le cash-flow ou flux de trésorerie est un indicateur financier qui mesure la capacité d’une entreprise à générer de la liquidité à partir de ses activités opérationnelles. Il représente la différence entre les entrées et sorties d’argent sur une période donnée.Elle reflète la capacité de l’entreprise à générer de la trésorerie à partir de son activité. Pour calculer dont notre cash-flow nous allons faire usage de la formule suivant.


\textbf{Formules d’usage :}
\begin{itemize}
    \item Cash-flow = Résultat net + (dotation nette aux amortissements et aux provisions) $-$ (plus-values de cession d’actifs) + (moins-values de cession d’actifs) $-$ (variation du besoin en fonds de roulement)
    \item Cash-flow actualisé = Cash-flow $\times (1 + t)^{-n}$ avec $t = 10\,\%$ et $n = \text{nombre d'années}$
\end{itemize}

\begin{xltabular}{\textwidth}{@{} X r r r @{}}
\caption{Cash-flow prévisionnel sur trois années} \label{tab:cash_flow} \\
\toprule
\rowcolor{BleuNuit}\color{white}\textbf{Éléments} & \rowcolor{BleuNuit}\color{white}\textbf{Année N} & \rowcolor{BleuNuit}\color{white}\textbf{Année N+1} & \rowcolor{BleuNuit}\color{white}\textbf{Année N+2} \\
\midrule
\endfirsthead
\multicolumn{4}{c}{\tablename\ \thetable{} -- Suite de la page précédente} \\
\toprule
\rowcolor{BleuNuit}\color{white}\textbf{Éléments} & \rowcolor{BleuNuit}\color{white}\textbf{Année N} & \rowcolor{BleuNuit}\color{white}\textbf{Année N+1} & \rowcolor{BleuNuit}\color{white}\textbf{Année N+2} \\
\midrule
\endhead
\midrule
\multicolumn{4}{r}{\textit{Suite à la page suivante...}} \\
\endfoot
\bottomrule
\multicolumn{4}{l}{\textit{Source : Les promoteurs 2025-2026}} \\
\endlastfoot
Résultat net & 1 867 916 & 30 545 335 & 76 165 640 \\
\rowcolor{BleuClair}Dotation aux amortissement & 1 957 340 & 1 957 340 & 1 957 340 \\
\midrule
\rowcolor{BleuClair}\textbf{Cash-flow} & \textbf{3 825 256} & \textbf{32 502 675} & \textbf{78 122 980} \\
Cash-flow actualisé & 3 477 505 & 29 547 886 & 71 020 891 \\
Cumul des cash-flows actualisés & 3 477 505 & 33 025 392 & 104 046 283 \\
\end{xltabular}

\subsubsection{Valeur actualisée nette (VAN)}
La VAN mesure l’avantage absolu susceptible d’être retiré d’un projet. Elle est égale à la différence entre les flux nets de trésorerie actualisés sur la durée de vie de l’investissement et le montant du capital initial investi.

La VAN peut être calculée directement avec le tableur Excel. Le taux d’actualisation retenu est de 10\,\%. Nous conservons la série de valeurs comprenant l’investissement initial et les cash-flows issus du compte de résultat entre l’année 1 et l’année 3, tels qu’indiqués ci-dessus.

$$\text{VAN} = -I_0 + \sum \text{CF} \cdot (1+t)^{-n}$$
avec :
\begin{itemize}
    \item $I_0$ = investissement initial ;
    \item $n$ = nombre d’années ;
    \item $t$ = taux d’actualisation du capital ;
    \item $\text{CF}$ = Cash-flow ;
    \item Cash-flow net = résultat net + dotations aux amortissements ;
    \item Cash-flow actualisé = Cash-flow $(1 + 10\%)^{-\text{n}}$ ;
    \item Si la VAN est positive, alors le projet est rentable ;
    \item Si la VAN est négative, alors le projet n’est pas suffisamment rentable.
\end{itemize}

\textbf{Application numérique :} la VAN obtenue au taux d’actualisation de 10\,\% est de \textbf{55\ 572\ 480 FCFA}, conformément au Tableau~\ref{tab:van}.

Le caractère strictement positif de la VAN (55\ 572\ 480 FCFA) valide l'hypothèse de rentabilité du projet. Les flux de trésorerie actualisés couvrent l'investissement initial et créent une valeur nette positive sur la période d'exploitation.

\begin{xltabular}{\textwidth}{@{} X r r r r @{}}
\caption{Valeur actualisée nette} \label{tab:van} \\
\toprule
\rowcolor{BleuNuit}\color{white}\textbf{Éléments} & \rowcolor{BleuNuit}\color{white}\textbf{Montants} & \rowcolor{BleuNuit}\color{white}\textbf{Année N} & \rowcolor{BleuNuit}\color{white}\textbf{Année N+1} & \rowcolor{BleuNuit}\color{white}\textbf{Année N+2} \\
\midrule
\endfirsthead
\multicolumn{5}{c}{\tablename\ \thetable{} -- Suite de la page précédente} \\
\toprule
\rowcolor{BleuNuit}\color{white}\textbf{Éléments} & \rowcolor{BleuNuit}\color{white}\textbf{Montants} & \rowcolor{BleuNuit}\color{white}\textbf{Année N} & \rowcolor{BleuNuit}\color{white}\textbf{Année N+1} & \rowcolor{BleuNuit}\color{white}\textbf{Année N+2} \\
\midrule
\endhead
\midrule
\multicolumn{5}{r}{\textit{Suite à la page suivante...}} \\
\endfoot
\bottomrule
\multicolumn{5}{l}{\textit{Source : Les promoteurs 2025-2026}} \\
\endlastfoot
Cumuls des cash-flows nets & -22 550 500 & 3 825 256 & 32 502 675 & 78 122 980 \\
\rowcolor{BleuClair}Taux & 10\,\% & & & \\
\textbf{Valeur actuelle nette (VAN)} & \textbf{55 572 480} & & & \\
\end{xltabular}
Le calcul de la VAN indique une valeur de 55\ 572\ 480 FCFA. Elle représente le surplus généré par les flux prévisionnels actualisés au taux de 10\,\% sur trois ans, après déduction de l'investissement initial.
\subsection{Taux de rentabilité interne (TRI)}
Le Taux de Rentabilité Interne (TRI) est un indicateur financier clé qui permet d'évaluer la rentabilité d'un projet en calculant le taux d'actualisation qui annule la valeur actuelle nette (VAN). Autrement dit, c'est le taux pour lequel les flux de trésorerie actualisés sont égaux à l'investissement initial.
\begin{itemize}
    \item \textbf{Formules}:
        - CF\_t = Cash-flow à l'année t
        - t = taux de rentabilité interne (TRI)
        - I\_0 = investissement initial
        - n = nombre d'années de prévision
        $I_0 = \sum \frac{CF_t}{(1 + i)^t}$
    \item \textbf{Critères:}
    - Si le TRI > taux d'actualisation (coût du capital) → le projet est rentable
    - Si le TRI < taux d'actualisation → le projet est non rentable
    - Si le TRI = taux d'actualisation → le projet est à l'équilibre (VAN = 0)
\end{itemize}

\subsection{Délai de récupération du capital investi (DRCI)}

Le délai de récupération correspond au délai au bout duquel le montant calculé des cashflow actualisés est égal au montant du capital investi. Il permet d'appréhender le délai nécessaire à la récupération du capital investi en le comparant aux flux cumulés de trésorerie.

\begin{itemize}
    \item Calcul du Délai de récupération du capital investi (DRCI) :
    \item $\text{DRCI} = \frac{I_0 - \text{cumul inférieur}}{\text{cumul supérieur} - \text{cumul inférieur}} + 1$
    \item $\text{DRCI} = \frac{22\ 550\ 500 - 3\ 477\ 505}{33\ 025\ 392 - 3\ 477\ 505} + 1 = 1{,}6455$

    \item $\text{DRCI} \approx 1{,}65$ an, soit environ \textbf{1 an, 7 mois et 23 jours}.
\end{itemize}

\begin{xltabular}{\textwidth}{@{} X r r r r @{}}
\caption{Calcul du TRI et DRCI} \label{tab:tri_drci} \\
\toprule
\rowcolor{BleuNuit}\color{white}\textbf{Éléments} & \rowcolor{BleuNuit}\color{white}\textbf{Montants} & \rowcolor{BleuNuit}\color{white}\textbf{Année N} & \rowcolor{BleuNuit}\color{white}\textbf{Année N+1} & \rowcolor{BleuNuit}\color{white}\textbf{Année N+2} \\
\midrule
\endfirsthead
\multicolumn{5}{c}{\tablename\ \thetable{} -- Suite de la page précédente} \\
\toprule
\rowcolor{BleuNuit}\color{white}\textbf{Éléments} & \rowcolor{BleuNuit}\color{white}\textbf{Montants} & \rowcolor{BleuNuit}\color{white}\textbf{Année N} & \rowcolor{BleuNuit}\color{white}\textbf{Année N+1} & \rowcolor{BleuNuit}\color{white}\textbf{Année N+2} \\
\midrule
\endhead
\midrule
\multicolumn{5}{r}{\textit{Suite à la page suivante...}} \\
\endfoot
\bottomrule
\multicolumn{5}{l}{\textit{Source : Les promoteurs 2025-2026}} \\
\endlastfoot
Investissement initial & -22 550 500 & & & -22 550 500 \\
\rowcolor{BleuClair}Cash-flow bruts & & -89 424 & 28 587 995 & 74 208 300 \\
Cash-flow actualisés (10\,\%) & & 3 477 505 & 29 547 886 & 71 020 891 \\
\midrule
Cash-flows cumulés & -22 550 500 & 3 477 505 & 33 025 392 & 104 046 283 \\
\midrule
\textbf{TRI} & \textbf{89,5\,\%} & & & \\
\end{xltabular}

Notre capital investi est totalement récupéré au bout d'environ 1 an, 7 mois et 23 jours. L'indice de profitabilité (IP) est de 3,46.

\subsubsection{Indice de profitabilité (IP)}
L’indice de profitabilité est le rapport entre le cumul actualisé des flux de trésorerie et le capital investi. En d’autres termes, c’est le rapport entre la somme des cash-flows actualisés et l’investissement initial.
$
IP = VAN/Io + 1
IP =55 572 480 /22 550 500 + 1
IP = 3.46^3$
Celui-ci signifie que pour chaque franc investi, l’entreprise génère 3 francs de profit. En d’autres termes, cet indicateur montre une rentabilité très élevée, suggérant que notre entreprise est extrêmement profitable et qu’elle utilise efficacement ses ressources pour générer des gains substantiels.

\section{Bilan de la viabilité financière}

L’analyse financière du projet ARGOS, adossée au modèle innovant de partage de la valeur ajoutée et ancrée sur les flux réels des ports de la sous-région, démontre une robustesse et une viabilité remarquables à moyen terme. S'appuyant sur un marché en forte expansion matérialisé par un trafic conteneurisé en Afrique centrale dépassant les 2,19 millions d'EVP et une croissance annuelle moyenne de +22,3\,\% enregistrée à Kribi, le modèle économique valide une synergie gagnant-gagnant où ARGOS perçoit 60\,\% des revenus par trajet, tandis que l’armateur partenaire conserve 40\,\% pour couvrir la maintenance de premier niveau et les aléas opérationnels. 

Cette architecture financière se traduit par une rentabilité nette démontrée dès la deuxième année d'exploitation, dégageant un résultat positif de 28\,587\,995 FCFA, qui s'accélère pour atteindre 74\,208\,300 FCFA dès la troisième année. Les indicateurs d'évaluation confirment l'attractivité exceptionnelle de l'opportunité d'investissement, illustrée par une valeur actualisée nette (VAN) hautement positive, un indice de profitabilité (IP) de 3,46 et un délai de récupération du capital investi (DRCI) d'environ un an et huit mois. Même en simulant un scénario conservateur fondé sur seulement quatre rotations annuelles par dispositif, la viabilité du projet demeure intacte. Portée par une marge brute de premier plan, l'atteinte du seuil de rentabilité dès le second exercice et un chiffre d'affaires multiplié par quatre entre la première et la troisième année, la trajectoire financière d'ARGOS confirme la solidité de sa proposition de valeur sur le marché de la logistique intelligente.


\indent Le projet ARGOS présente une viabilité financière solide et constitue une opportunité d’investissement attractive. La stratégie de déploiement progressif et le partage des bénéfices avec les entreprises partenaires garantissent une adoption fluide et une pérennité du modèle économique.
